\chapter{ارزیابی و تحلیل نتایج}

\section{مقدمه}
\label{sec:evaluation_intro}

در فصل‌های پیشین، ابتدا مبانی نظری و پیشینه پژوهش در زمینه ناوبری خودکار، یادگیری تقویتی، ایمنی تصمیم‌گیری و استفاده از سیاست‌های پیوسته بررسی شد. سپس، چارچوب پیشنهادی این پژوهش برای ناوبری اسکوتر چهارچرخ در محیط‌های دینامیکی تشریح گردید. این چارچوب شامل طراحی فضای مشاهده ترکیبی، فضای عمل پیوسته، تابع پاداش چندهدفه، سپر ایمنی غیرفعال، معماری سیاست \lr{LSTM-SAC} و روند آموزش مرحله‌ای است. در ادامه نیز اجزای پیاده‌سازی عملی سامانه، شامل معماری سخت‌افزاری و نرم‌افزاری، دریافت داده از حسگرها، تولید مشاهده واقعی و اجرای مدل آموزش‌دیده روی سامانه عملی توضیح داده شد.

با توجه به ساختار روش پیشنهادی در فصل چهارم، فرآیند آموزش در این پژوهش تنها به آموزش یک عامل منفرد محدود نیست. ابتدا یک عامل معلم با دسترسی به هندسه دقیق شبیه‌سازی آموزش داده می‌شود تا یک سیاست مرجع نسبتاً پایدار در شرایط مشاهده تمیز به‌دست آید. سپس عامل دانشجو با استفاده از داده‌های ادراکی و با کمک لنگر رفتاری حاصل از سیاست معلم آموزش داده می‌شود. در این مرحله، هدف آن است که سیاست دانشجو بتواند در شرایطی نزدیک‌تر به پیاده‌سازی واقعی، یعنی در حضور خطاها و محدودیت‌های ادراکی، رفتار ناوبری ایمن و کارآمدی تولید کند. بنابراین، تحلیل نتایج باید میان نقش عامل معلم به‌عنوان سیاست مرجع و نقش عامل دانشجو به‌عنوان سیاست نهایی قابل انتقال به سامانه عملی تمایز قائل شود.

هدف فصل حاضر، ارزیابی و تحلیل نتایج حاصل از این روند آموزشی و اجرایی است. در این فصل بررسی می‌شود که آیا سیاست آموزش‌دیده، به‌ویژه سیاست دانشجو، توانسته است میان سه هدف اصلی ناوبری، یعنی رسیدن به هدف، حفظ ایمنی و تولید حرکت نرم و قابل اجرا، تعادل مناسبی برقرار کند یا خیر. ازاین‌رو، تحلیل نتایج تنها به مقدار پاداش محدود نمی‌شود، بلکه مجموعه‌ای از شاخص‌های مکمل شامل توزیع دلایل پایان اپیزود، طول اپیزود، روند زیان‌های شبکه، همواری حرکت، حداقل فاصله از موانع، وضعیت سپر ایمنی و اثر آموزش مرحله‌ای نیز بررسی می‌شوند.

تحلیل‌های این فصل در سه سطح انجام می‌گیرد. در سطح نخست، روند آموزش در محیط شبیه‌سازی بررسی می‌شود تا مشخص گردد فرآیند یادگیری عامل در طول مراحل مختلف آموزش چگونه شکل گرفته و با افزایش تدریجی دشواری سناریوها چه تغییری کرده است. در این بخش، تمرکز اصلی بر نتایج آموزش دانشجو است، زیرا این عامل بر پایه مشاهده ادراکی آموزش دیده و سیاست نهایی مورد استفاده در ارزیابی و پیاده‌سازی عملی را تشکیل می‌دهد. در صورت نیاز، نتایج عامل معلم نیز به‌عنوان مرجع آموزشی و برای تفسیر روند انتقال از هندسه دقیق به ادراک واقعی مورد استفاده قرار می‌گیرد.

در سطح دوم، سیاست نهایی در سناریوهای آزمون ارزیابی می‌شود. در این مرحله، پارامترهای شبکه ثابت بوده و عامل دیگر در حال یادگیری نیست؛ بنابراین، نتایج به‌دست‌آمده نشان‌دهنده رفتار نهایی سیاست در شرایط ارزیابی هستند. این تحلیل مشخص می‌کند که سیاست نهایی تا چه حد در رسیدن به هدف، اجتناب از برخورد، حفظ فاصله مناسب از موانع و تولید فرمان‌های قابل اجرا موفق بوده است.

در سطح سوم، نتایج مرتبط با پیاده‌سازی عملی سامانه مورد بحث قرار می‌گیرد. هدف این بخش، بررسی این موضوع است که اجزای ادراکی و اجرایی سامانه، از جمله فاصله‌سنجی استریو، حسگرهای نزدیک‌برد، زنجیره تولید مشاهده واقعی و اجرای مدل روی سامانه عملی، تا چه حد امکان انتقال روش طراحی‌شده از شبیه‌سازی به نمونه واقعی را فراهم کرده‌اند.

بنابراین، فصل حاضر نقش جمع‌بندی تحلیلی روش پیشنهادی را بر عهده دارد. نتایج این فصل نشان می‌دهند که روش پیشنهادی در چه بخش‌هایی عملکرد قابل دفاعی دارد، در کدام بخش‌ها با محدودیت مواجه است و تا چه اندازه می‌توان سیاست نهایی را برای ناوبری ایمن اسکوتر چهارچرخ قابل اتکا دانست.

\section{پروتکل ارزیابی و شاخص‌های عملکرد}
\label{sec:evaluation_protocol}

برای تحلیل عملکرد عامل‌های یادگیرنده، لازم است پیش از ارائه نتایج، چارچوب ارزیابی و شاخص‌های مورد استفاده به‌صورت روشن تعریف شوند. در مسئله ناوبری اسکوتر چهارچرخ، موفقیت عامل تنها با رسیدن به هدف سنجیده نمی‌شود؛ زیرا یک سیاست ممکن است با عبور از مسیرهای پرخطر، نزدیک‌شدن بیش از حد به موانع یا تولید فرمان‌های ناگهانی نیز به هدف برسد. چنین رفتاری، حتی اگر در شبیه‌سازی به موفقیت منجر شود، برای سامانه واقعی قابل قبول نیست. ازاین‌رو، ارزیابی در این پژوهش بر اساس مجموعه‌ای از شاخص‌های مکمل انجام می‌شود که سه جنبه اصلی عملکرد، یعنی کارایی ناوبری، ایمنی حرکت و کیفیت فرمان‌های کنترلی را پوشش می‌دهند.

با توجه به ساختار آموزش معلم--دانشجو، ارزیابی نتایج در این فصل در دو سطح آموزشی و یک سطح اجرایی انجام می‌شود. در سطح آموزشی، روند یادگیری عامل‌ها در طول اپیزودهای آموزش بررسی می‌گردد. در این بخش، عامل معلم نقش سیاست مرجع را دارد و عامل دانشجو به‌عنوان سیاست اصلی مورد توجه قرار می‌گیرد؛ زیرا دانشجو با داده‌های ادراکی و تحت شرایط نزدیک‌تر به پیاده‌سازی واقعی آموزش دیده است. در سطح آزمون نهایی، سیاست آموزش‌دیده بدون به‌روزرسانی پارامترهای شبکه در سناریوهای ارزیابی اجرا می‌شود. در سطح اجرایی نیز نتایج مرتبط با پیاده‌سازی عملی، به‌ویژه قابلیت تولید مشاهده واقعی و کیفیت داده‌های ادراکی، مورد بررسی قرار می‌گیرد.

\subsection{شرایط کلی ارزیابی}
\label{subsec:evaluation_conditions}

در مرحله آموزش، عامل‌ها در محیط شبیه‌سازی‌شده مبتنی بر \lr{Gazebo} و \lr{ROS2} با مدل اسکوتر چهارچرخ، هدف، موانع محیطی و حسگرهای مجازی تعامل می‌کنند. در هر اپیزود، عامل بر اساس مشاهده دریافتی از محیط، فرمان سرعت خطی و سرعت زاویه‌ای را تولید می‌کند و محیط در پاسخ به این فرمان، وضعیت بعدی، مقدار پاداش و وضعیت پایان اپیزود را تعیین می‌نماید. داده‌های مربوط به پاداش، دلیل پایان اپیزود، طول اپیزود، زیان‌های شبکه، وضعیت سپر ایمنی، فاصله از موانع و شاخص‌های نرمی حرکت در طول فرآیند آموزش ثبت می‌شوند.

در آموزش معلم، مشاهده عامل بر پایه هندسه دقیق شبیه‌سازی ساخته می‌شود. هدف از این مرحله، به‌دست‌آوردن سیاستی است که در شرایط مشاهده دقیق بتواند رفتار ناوبری پایداری تولید کند. در مقابل، در آموزش دانشجو، عامل با مشاهده ادراکی و با استفاده از لنگر رفتاری معلم آموزش می‌بیند. بنابراین، نتایج آموزش دانشجو از نظر انتقال به سامانه واقعی اهمیت بیشتری دارد، زیرا این عامل باید با داده‌هایی تصمیم‌گیری کند که به داده‌های قابل دریافت در پیاده‌سازی عملی نزدیک‌تر هستند.

در مرحله ارزیابی نهایی، سیاست آموزش‌دیده بدون به‌روزرسانی پارامترهای شبکه اجرا می‌شود. هدف از این مرحله، بررسی عملکرد سیاست نهایی در شرایطی است که عامل دیگر در حال یادگیری نیست و تنها بر اساس رفتار آموخته‌شده تصمیم‌گیری می‌کند. در این مرحله، سناریوهای آزمون از نظر تعداد موانع، حرکت موانع، حرکت هدف، موقعیت اولیه و پیچیدگی محیط طراحی می‌شوند تا قابلیت سیاست در شرایط مختلف بررسی گردد.

برای کاهش اثر تصادفی‌بودن شرایط اولیه، هر سناریو در چندین اپیزود مستقل اجرا می‌شود. در هر اجرا، وضعیت‌هایی مانند رسیدن به هدف، برخورد، اتمام زمان، طول اپیزود، حداقل فاصله از موانع، زمان تا برخورد، رخدادهای سپر ایمنی و شاخص‌های کیفیت فرمان ثبت می‌شوند. سپس شاخص‌های نهایی با استفاده از میانگین‌گیری، محاسبه نسبت رخدادها و بررسی روندهای آماری تحلیل می‌گردند. این رویکرد باعث می‌شود قضاوت درباره عملکرد عامل بر اساس یک اجرای منفرد یا یک سناریوی خاص انجام نشود.

\subsection{شاخص‌های ارزیابی}
\label{subsec:evaluation_metrics}

شاخص‌های مورد استفاده در این فصل به‌گونه‌ای انتخاب شده‌اند که بتوانند رفتار عامل را از چند دیدگاه مختلف بررسی کنند. در تحلیل آموزش، شاخص‌هایی مانند پاداش تجمعی، توزیع دلایل پایان اپیزود، طول اپیزود، زیان‌های شبکه، شاخص‌های نرمی حرکت، توزیع عملکرد در مراحل مختلف آموزش و وضعیت سپر ایمنی بررسی می‌شوند. این شاخص‌ها با نمودارهای منتخب آموزش شامل روند توزیعی پاداش، توزیع دلایل پایان اپیزود، طول اپیزود، زیان کنشگر و منتقد، شاخص \lr{Jerk}، خلاصه مرحله‌ای شاخص‌ها، ترکیب نتایج در مراحل آموزش، وابستگی به سپر، شدت رخدادهای سپر و حداقل فاصله از موانع تحلیل می‌گردند.

نرخ موفقیت نشان‌دهنده نسبت اپیزودهایی است که در آن‌ها عامل به هدف رسیده است. اگر تعداد کل اپیزودهای بررسی‌شده در یک بازه برابر با \(N\) و تعداد اپیزودهای موفق برابر با \(N_{\mathrm{goal}}\) باشد، نرخ موفقیت به‌صورت زیر تعریف می‌شود:
\begin{equation}
	SR = \frac{N_{\mathrm{goal}}}{N}
	\label{eq:success_rate_definition}
\end{equation}

در کنار نرخ موفقیت، نرخ برخورد و نرخ اتمام زمان نیز بررسی می‌شوند. اگر تعداد اپیزودهای منتهی به برخورد برابر با \(N_{\mathrm{collision}}\) و تعداد اپیزودهای منتهی به اتمام زمان برابر با \(N_{\mathrm{timeout}}\) باشد، این دو شاخص به‌صورت زیر محاسبه می‌شوند:
\begin{equation}
	CR = \frac{N_{\mathrm{collision}}}{N},
	\qquad
	TR = \frac{N_{\mathrm{timeout}}}{N}
	\label{eq:collision_timeout_rate_definition}
\end{equation}

پاداش تجمعی هر اپیزود برابر مجموع پاداش‌ها و جریمه‌هایی است که عامل در طول آن اپیزود دریافت می‌کند. اگر پاداش لحظه‌ای در گام زمانی \(t\) برابر با \(r_t\) و طول اپیزود برابر با \(T\) باشد، پاداش تجمعی اپیزود به‌صورت زیر تعریف می‌شود:
\begin{equation}
	R_{\mathrm{ep}} = \sum_{t=0}^{T} r_t
	\label{eq:episode_return_definition}
\end{equation}

با توجه به چندمؤلفه‌ای بودن تابع پاداش، افزایش پاداش به‌تنهایی برای قضاوت درباره کیفیت سیاست کافی نیست. بنابراین، روند پاداش باید همراه با نرخ موفقیت، دلایل پایان اپیزود، برخورد، اتمام زمان، وضعیت سپر ایمنی و فاصله از موانع تفسیر شود. به همین دلیل، در تحلیل آموزش از نمودار توزیعی پاداش استفاده می‌شود تا علاوه بر روند مرکزی، پراکندگی و نوسان رفتار عامل نیز قابل مشاهده باشد.

طول اپیزود برابر تعداد گام‌های زمانی است که عامل تا رسیدن به یکی از شرایط خاتمه طی می‌کند. کاهش طول اپیزود زمانی مطلوب است که اپیزود با رسیدن به هدف پایان یافته باشد. در مقابل، اگر کوتاه شدن اپیزود ناشی از برخورد زودهنگام باشد، نشانه بهبود نیست. بنابراین، طول اپیزود در این فصل همراه با دلیل پایان اپیزود تحلیل می‌شود.

توزیع دلایل پایان اپیزود یکی از شاخص‌های کلیدی برای تشخیص ماهیت واقعی یادگیری است. این شاخص نشان می‌دهد که در طول آموزش، چه نسبتی از اپیزودها با رسیدن به هدف، برخورد یا اتمام زمان پایان یافته‌اند. اگر افزایش پاداش همراه با افزایش سهم اپیزودهای موفق و کاهش سهم برخورد و اتمام زمان باشد، می‌توان آن را نشانه‌ای از بهبود واقعی سیاست دانست. در غیر این صورت، افزایش پاداش ممکن است ناشی از رفتارهایی باشد که از نظر ناوبری یا ایمنی قابل قبول نیستند.

زیان کنشگر و منتقد برای بررسی پایداری فرآیند یادگیری استفاده می‌شوند. زیان منتقد به خطای تخمین تابع ارزش مربوط است و زیان کنشگر روند اصلاح سیاست را نشان می‌دهد. نوسان شدید یا رشد ناپایدار این زیان‌ها می‌تواند نشانه‌ای از بی‌ثباتی آموزش، نامناسب بودن مقیاس پاداش یا دشواری بیش از حد محیط باشد. بااین‌حال، مقدار مطلق زیان‌ها به‌تنهایی شاخص مستقیمی از کیفیت ناوبری نیست و باید همراه با پاداش و شاخص‌های عملکردی تفسیر شود.

شاخص‌های ایمنی شامل نرخ برخورد، حداقل فاصله از موانع، وضعیت سپر ایمنی و شدت رخدادهای سپر هستند. نرخ برخورد تنها وقوع برخورد نهایی را نشان می‌دهد، اما وضعیت سپر ایمنی می‌تواند رفتارهای پرخطر قبل از برخورد را نیز آشکار کند. بنابراین، اگر برخورد رخ ندهد ولی عامل بارها وارد ناحیه خطر یا بحرانی شود، نمی‌توان سیاست را کاملاً ایمن دانست. به همین دلیل، در این فصل علاوه بر نرخ برخورد، روند وابستگی به سپر و شدت رخدادهای هشدار، خطر و بحرانی نیز تحلیل می‌شود.

حداقل فاصله از موانع نشان می‌دهد عامل در طول اپیزود تا چه اندازه به نزدیک‌ترین مانع نزدیک شده است. اگر فاصله عامل از مانع \(i\) در گام زمانی \(t\) برابر با \(d_i(t)\) باشد، حداقل فاصله اپیزود به‌صورت زیر نوشته می‌شود:
\begin{equation}
	d_{\min} = \min_{t,i} d_i(t).
	\label{eq:min_obstacle_distance_definition}
\end{equation}

کیفیت حرکت نیز برای قابلیت اجرای فرمان‌ها روی سامانه واقعی اهمیت دارد. سیاستی که فرمان‌های ناگهانی یا پرنوسان تولید کند، حتی در صورت موفقیت در شبیه‌سازی، ممکن است در نمونه واقعی عملکرد مناسبی نداشته باشد. برای بررسی این موضوع، تغییرات فرمان سرعت خطی و سرعت زاویه‌ای و شاخص تکان \lr{Jerk} تحلیل می‌شوند. اگر سرعت خطی و زاویه‌ای در دو گام زمانی متوالی به‌ترتیب برابر با \(v_t\)، \(v_{t-1}\)، \(\omega_t\) و \(\omega_{t-1}\) باشند، تغییرات فرمان به‌صورت زیر تعریف می‌شود:
\begin{equation}
	\Delta v_t = v_t - v_{t-1},
	\qquad
	\Delta \omega_t = \omega_t - \omega_{t-1}.
	\label{eq:control_variation_definition}
\end{equation}

برای بررسی کلی نوسان فرمان در طول اپیزود، می‌توان مجموع مربع تغییرات فرمان را به‌صورت زیر در نظر گرفت:
\begin{equation}
	J_{\mathrm{ctrl}} =
	\sum_{t=1}^{T}
	\left[
	(\Delta v_t)^2 +
	(\Delta \omega_t)^2
	\right].
	\label{eq:control_smoothness_definition}
\end{equation}

کاهش این شاخص نشان می‌دهد سیاست از فرمان‌های ناگهانی و پرنوسان اجتناب کرده است. البته نرمی بیش از حد نیز همیشه مطلوب نیست، زیرا ممکن است واکنش عامل به موانع متحرک را کند کند. بنابراین، شاخص‌های کیفیت حرکت باید همراه با نرخ موفقیت، فاصله از موانع و وضعیت سپر ایمنی تحلیل شوند.

\section{تحلیل نتایج آموزش در محیط شبیه‌سازی}
\label{sec:training_results_analysis}

در این بخش، نتایج آموزش عامل معلم و عامل دانشجو در محیط شبیه‌سازی تحلیل می‌شود. هدف اصلی، بررسی این موضوع است که آیا روند آموزش معلم--دانشجو توانسته است سیاستی با عملکرد قابل قبول، پایداری نسبی، فرمان‌های قابل اجرا و رفتار ایمن‌تر برای عامل دانشجو ایجاد کند یا خیر. در تحلیل حاضر، عامل معلم به‌عنوان سیاست مرجع مبتنی بر مشاهده هندسی دقیق و عامل دانشجو به‌عنوان سیاست نهایی مبتنی بر مشاهده ادراکی در نظر گرفته می‌شود.

\subsection{روند کلی یادگیری}
\label{subsec:overall_training_trend}

شکل~\ref{fig:teacher_student_reward_distribution} روند توزیع پاداش تجمعی را در طول آموزش برای عامل معلم و عامل دانشجو نشان می‌دهد. در هر دو عامل، پاداش در ابتدای آموزش از مقادیر منفی شروع شده و پس از تعداد کمی اپیزود به‌سرعت وارد ناحیه مثبت می‌شود. این جهش اولیه نشان می‌دهد که هر دو سیاست توانسته‌اند از رفتارهای کاملاً ناموفق و اکتشافی فاصله بگیرند و الگوی پایه‌ای حرکت به سمت هدف را یاد بگیرند.

\begin{figure}[htbp]
	\centering
	
	\subfloat[عامل معلم]{
		\begin{minipage}{\textwidth}
			\centering
			\includegraphics[width=0.95\textwidth, height=5cm]{17_reward_quantile_band_teacher.png}
			
			\vspace{0.4em}
			
			{\small
				\href{https://mohammadabedi1179.github.io/thesis-archive/chapters/ch6/17_reward_quantile_band_teacher.html}{\textbf{مشاهده نسخه تعاملی شکل با کیفیت بالاتر}}
				\qquad
				\qrcode[height=1.2cm]{https://mohammadabedi1179.github.io/thesis-archive/chapters/ch6/17_reward_quantile_band_teacher.html}
			}
		\end{minipage}
	}
	
	\vspace{0.8em}
	
	\subfloat[عامل دانشجو]{
		\begin{minipage}{\textwidth}
			\centering
			\includegraphics[width=0.95\textwidth, height=5cm]{17_reward_quantile_band_student.png}
			
			\vspace{0.4em}
			
			{\small
				\href{https://mohammadabedi1179.github.io/thesis-archive/chapters/ch6/17_reward_quantile_band_student.html}{\textbf{مشاهده نسخه تعاملی شکل با کیفیت بالاتر}}
				\qquad
				\qrcode[height=1.2cm]{https://mohammadabedi1179.github.io/thesis-archive/chapters/ch6/17_reward_quantile_band_student.html}
			}
		\end{minipage}
	}
	
	\caption{مقایسه روند توزیع پاداش در طول آموزش: الف) عامل معلم، ب) عامل دانشجو}
	\label{fig:teacher_student_reward_distribution}
\end{figure}

پس از فاز ابتدایی، مقدار مرکزی پاداش در هر دو نمودار در سطح نسبتاً بالایی باقی می‌ماند؛ بااین‌حال، نوسان‌ها و نقاط پرت منفی نشان می‌دهند که عملکرد عامل‌ها در تمام اپیزودها یکنواخت نیست. این موضوع به‌ویژه برای عامل دانشجو مهم‌تر است، زیرا مشاهده ادراکی نسبت به مشاهده هندسی دقیق عدم‌قطعیت بیشتری دارد. بنابراین، روند پاداش نشان‌دهنده شکل‌گیری رفتار مفید است، اما برای قضاوت نهایی باید همراه با دلیل پایان اپیزود و طول اپیزود تفسیر شود.

شکل~\ref{fig:teacher_student_termination_reason_distribution} نشان می‌دهد که در بخش عمده آموزش، اپیزودهای موفق سهم غالب را تشکیل می‌دهند. در عامل معلم، برخوردها محدودتر و پراکنده‌تر هستند، در حالی که در عامل دانشجو، سهم برخوردها در مراحل میانی و پایانی افزایش بیشتری دارد. این نتیجه نشان می‌دهد که آموزش دانشجو موفق بوده، اما انتقال از مشاهده هندسی به مشاهده ادراکی با کاهش نسبی حاشیه ایمنی همراه شده است.

\begin{figure}[htbp]
	\centering
	
	\subfloat[]{
		\begin{minipage}{0.75\textwidth}
			\centering
			
			\includegraphics[width=0.9\textwidth]
			{04_reason_distribution_teacher.png}
			
			\vspace{0.35em}
			
			{\small
				\href{https://mohammadabedi1179.github.io/thesis-archive/chapters/ch6/04_reason_distribution_teacher.html}
				{\textbf{مشاهده نسخه تعاملی عامل معلم با کیفیت بالاتر}}
				\hspace{0.8em}
				\raisebox{-0.4\height}{
					\qrcode[height=1.2cm]{https://mohammadabedi1179.github.io/thesis-archive/chapters/ch6/04_reason_distribution_teacher.html}
				}
			}
		\end{minipage}
	}
	
	\vspace{0.7em}
	
	\subfloat[]{
		\begin{minipage}{0.75\textwidth}
			\centering
			
			\includegraphics[width=0.9\textwidth]
			{04_reason_distribution_student.png}
			
			\vspace{0.35em}
			
			{\small
				\href{https://mohammadabedi1179.github.io/thesis-archive/chapters/ch6/04_reason_distribution_student.html}
				{\textbf{مشاهده نسخه تعاملی عامل دانشجو با کیفیت بالاتر}}
				\hspace{0.8em}
				\raisebox{-0.4\height}{
					\qrcode[height=1.2cm]{https://mohammadabedi1179.github.io/thesis-archive/chapters/ch6/04_reason_distribution_student.html}
				}
			}
		\end{minipage}
	}
	
	\caption{مقایسه توزیع دلایل پایان اپیزود در طول آموزش:
		الف) عامل معلم، ب) عامل دانشجو}
	\label{fig:teacher_student_termination_reason_distribution}
\end{figure}

تحلیل طول اپیزود در شکل~\ref{fig:teacher_student_episode_length} همین برداشت را تکمیل می‌کند. در هر دو عامل، طول اپیزود پس از چند اپیزود ابتدایی به‌سرعت کاهش می‌یابد و در ادامه در بازه‌ای محدودتر قرار می‌گیرد. بااین‌حال، پراکندگی طول اپیزود در عامل دانشجو بیشتر است و بخشی از اپیزودهای کوتاه آن به برخورد منتهی شده‌اند. بنابراین، کاهش طول اپیزود برای دانشجو را نباید صرفاً نشانه کارایی بیشتر دانست؛ بخشی از آن ناشی از پایان زودهنگام نامطلوب است.

\begin{figure}[htbp]
	\centering
	
	\subfloat[]{
		\begin{minipage}{0.75\textwidth}
			\centering
			
			\includegraphics[width=0.9\textwidth]
			{03_episode_length_time_teacher.png}
			
			\vspace{0.35em}
			
			{\small
				\href{
					https://mohammadabedi1179.github.io/thesis-archive/chapters/ch6/03_episode_length_time_teacher.html
				}{
					\textbf{مشاهده نسخه تعاملی عامل معلم با کیفیت بالاتر}
				}
				\hspace{0.8em}
				\raisebox{-0.4\height}{
					\qrcode[height=1.2cm]{
						https://mohammadabedi1179.github.io/thesis-archive/chapters/ch6/03_episode_length_time_teacher.html
					}
				}
			}
			
		\end{minipage}
	}
	
	\vspace{0.7em}
	
	\subfloat[]{
		\begin{minipage}{0.75\textwidth}
			\centering
			
			\includegraphics[width=0.9\textwidth]
			{03_episode_length_time_student.png}
			
			\vspace{0.35em}
			
			{\small
				\href{
					https://mohammadabedi1179.github.io/thesis-archive/chapters/ch6/03_episode_length_time_student.html
				}{
					\textbf{مشاهده نسخه تعاملی عامل دانشجو با کیفیت بالاتر}
				}
				\hspace{0.8em}
				\raisebox{-0.4\height}{
					\qrcode[height=1.2cm]{
						https://mohammadabedi1179.github.io/thesis-archive/chapters/ch6/03_episode_length_time_student.html
					}
				}
			}
			
		\end{minipage}
	}
	
	\caption{مقایسه روند طول اپیزود در طول آموزش:
		الف) عامل معلم، ب) عامل دانشجو}
	\label{fig:teacher_student_episode_length}
\end{figure}



در مجموع، سه نمودار پاداش، دلیل پایان و طول اپیزود نشان می‌دهند که هر دو عامل به رفتار هدفمند رسیده‌اند، اما عامل معلم پایدارتر و ایمن‌تر عمل کرده است. عامل دانشجو نیز عملکرد کلی قابل قبول دارد، ولی در سناریوهای دشوارتر، نوسان و برخورد بیشتری نشان می‌دهد.

\subsection{پایداری یادگیری و کیفیت فرمان}
\label{subsec:learning_stability_motion_quality}

برای بررسی پایداری عددی فرآیند آموزش، شکل~\ref{fig:teacher_student_actor_critic_loss} روند زیان کنشگر و منتقد را نشان می‌دهد. در عامل معلم، زیان‌ها پس از فاز ابتدایی در محدوده‌ای کنترل‌شده باقی می‌مانند و نشانه‌ای از واگرایی جدی دیده نمی‌شود. در عامل دانشجو، زیان منتقد بزرگ‌تر و پراکنده‌تر است که با دشوارتر بودن تخمین ارزش در فضای مشاهده ادراکی سازگار است. با وجود این، روند کلی زیان‌ها به فروپاشی آموزش منجر نشده است.

\begin{figure}[htbp]
	\centering
	
	\subfloat[]{
		\begin{minipage}[t]{0.48\textwidth}
			\centering
			
			\includegraphics[width=\linewidth]
			{05_losses_entropy_teacher.png}
			
			\vspace{0.4em}
			
			{\small
				\href{https://mohammadabedi1179.github.io/thesis-archive/chapters/ch6/05_losses_entropy_teacher.html}
				{\textbf{مشاهده نسخه تعاملی عامل معلم با کیفیت بالاتر}}
				
				\par\vspace{0.3em}
				
				\href{https://mohammadabedi1179.github.io/thesis-archive/chapters/ch6/05_losses_entropy_teacher.html}{
					\qrcode[height=1.15cm]{https://mohammadabedi1179.github.io/thesis-archive/chapters/ch6/05_losses_entropy_teacher.html}
				}
			}
		\end{minipage}
	}
	\hfill
	\subfloat[]{
		\begin{minipage}[t]{0.48\textwidth}
			\centering
			
			\includegraphics[width=\linewidth]
			{05_losses_entropy_student.png}
			
			\vspace{0.4em}
			
			{\small
				\href{https://mohammadabedi1179.github.io/thesis-archive/chapters/ch6/05_losses_entropy_student.html}
				{\textbf{مشاهده نسخه تعاملی عامل دانشجو با کیفیت بالاتر}}
				
				\par\vspace{0.3em}
				
				\href{https://mohammadabedi1179.github.io/thesis-archive/chapters/ch6/05_losses_entropy_student.html}{
					\qrcode[height=1.15cm]{https://mohammadabedi1179.github.io/thesis-archive/chapters/ch6/05_losses_entropy_student.html}
				}
			}
		\end{minipage}
	}
	
	\caption{مقایسه روند زیان کنشگر و منتقد در طول آموزش:
		الف) عامل معلم، ب) عامل دانشجو}
	\label{fig:teacher_student_actor_critic_loss}
\end{figure}

شکل~\ref{fig:teacher_student_mean_jerk} کیفیت فرمان‌های تولیدشده را از نظر شاخص تکان میانگین نشان می‌دهد. در هر دو عامل، مقدار تکان در ابتدای آموزش بالا است و سپس کاهش می‌یابد. این کاهش نشان می‌دهد که سیاست‌ها به‌تدریج از تولید فرمان‌های ناگهانی فاصله گرفته‌اند. در هر دو نمودار، مؤلفه زاویه‌ای نسبت به مؤلفه خطی مقدار بیشتری دارد که با ماهیت مسئله ناوبری در حضور مانع و نیاز به اصلاح جهت سازگار است.

\begin{figure}[htbp]
	\centering
	
	\subfloat[]{
		\begin{minipage}{0.75\textwidth}
			\centering
			
			\includegraphics[width=0.9\textwidth]
			{08_mean_jerk_teacher.png}
			
			\vspace{0.35em}
			
			{\small
				\href{https://mohammadabedi1179.github.io/thesis-archive/chapters/ch6/08_mean_jerk_teacher.html}
				{\textbf{مشاهده نسخه تعاملی عامل معلم با کیفیت بالاتر}}
				\hspace{0.8em}
				\raisebox{-0.4\height}{
					\href{https://mohammadabedi1179.github.io/thesis-archive/chapters/ch6/08_mean_jerk_teacher.html}{
						\qrcode[height=1.2cm]{https://mohammadabedi1179.github.io/thesis-archive/chapters/ch6/08_mean_jerk_teacher.html}
					}
				}
			}
		\end{minipage}
	}
	
	\vspace{0.7em}
	
	\subfloat[]{
		\begin{minipage}{0.75\textwidth}
			\centering
			
			\includegraphics[width=0.9\textwidth]
			{08_mean_jerk_student.png}
			
			\vspace{0.35em}
			
			{\small
				\href{https://mohammadabedi1179.github.io/thesis-archive/chapters/ch6/08_mean_jerk_student.html}
				{\textbf{مشاهده نسخه تعاملی عامل دانشجو با کیفیت بالاتر}}
				\hspace{0.8em}
				\raisebox{-0.4\height}{
					\href{https://mohammadabedi1179.github.io/thesis-archive/chapters/ch6/08_mean_jerk_student.html}{
						\qrcode[height=1.2cm]{https://mohammadabedi1179.github.io/thesis-archive/chapters/ch6/08_mean_jerk_student.html}
					}
				}
			}
		\end{minipage}
	}
	
	\caption{مقایسه روند شاخص تکان میانگین در طول آموزش:
		الف) عامل معلم، ب) عامل دانشجو}
	\label{fig:teacher_student_mean_jerk}
\end{figure}

برآیند این دو نمودار نشان می‌دهد که آموزش تنها به افزایش پاداش محدود نشده، بلکه از نظر پایداری عددی و کیفیت فرمان نیز به رفتار قابل قبول رسیده است. بااین‌حال، زیان منتقد بزرگ‌تر در عامل دانشجو نشان می‌دهد که سیاست ادراکی نسبت به معلم با مسئله یادگیری دشوارتری روبه‌رو بوده است.

\subsection{اثر آموزش مرحله‌ای}
\label{subsec:curriculum_effect_analysis}

برای بررسی اثر مراحل مختلف آموزش، شکل~\ref{fig:teacher_student_stagewise_outcomes} ترکیب دلایل پایان اپیزود را به تفکیک stage نشان می‌دهد. در هر دو عامل، سهم اپیزودهای موفق در بیشتر مراحل غالب است؛ بنابراین، افزایش تدریجی دشواری سناریوها باعث فروپاشی عملکرد نشده است. بااین‌حال، در مراحل دشوارتر، سهم برخورد افزایش می‌یابد و این افزایش برای عامل دانشجو محسوس‌تر است.

\begin{figure}[htbp]
	\centering
	
	\subfloat[]{
		\begin{minipage}{0.75\textwidth}
			\centering
			
			\includegraphics[width=0.7\textwidth]
			{19_stagewise_outcome_composition_teacher.png}
			
			\vspace{0.35em}
			
			{\small
				\href{https://mohammadabedi1179.github.io/thesis-archive/chapters/ch6/19_stagewise_outcome_composition_teacher.html}
				{\textbf{مشاهده نسخه تعاملی عامل معلم با کیفیت بالاتر}}
				\hspace{0.8em}
				\raisebox{-0.4\height}{
					\href{https://mohammadabedi1179.github.io/thesis-archive/chapters/ch6/19_stagewise_outcome_composition_teacher.html}{
						\qrcode[height=1.2cm]
						{https://mohammadabedi1179.github.io/thesis-archive/chapters/ch6/19_stagewise_outcome_composition_teacher.html}
					}
				}
			}
			
		\end{minipage}
	}
	
	\vspace{0.7em}
	
	\subfloat[]{
		\begin{minipage}{0.75\textwidth}
			\centering
			
			\includegraphics[width=0.7\textwidth]
			{19_stagewise_outcome_composition_student.png}
			
			\vspace{0.35em}
			
			{\small
				\href{https://mohammadabedi1179.github.io/thesis-archive/chapters/ch6/19_stagewise_outcome_composition_student.html}
				{\textbf{مشاهده نسخه تعاملی عامل دانشجو با کیفیت بالاتر}}
				\hspace{0.8em}
				\raisebox{-0.4\height}{
					\href{https://mohammadabedi1179.github.io/thesis-archive/chapters/ch6/19_stagewise_outcome_composition_student.html}{
						\qrcode[height=1.2cm]
						{https://mohammadabedi1179.github.io/thesis-archive/chapters/ch6/19_stagewise_outcome_composition_student.html}
					}
				}
			}
			
		\end{minipage}
	}
	
	\caption{ترکیب دلایل پایان اپیزود به تفکیک مراحل مختلف آموزش:
		الف) عامل معلم، ب) عامل دانشجو}
	\label{fig:teacher_student_stagewise_outcomes}
\end{figure}

این نتیجه نشان می‌دهد که curriculum توانسته است یادگیری را در طول مراحل مختلف حفظ کند، اما افزایش دشواری محیط مستقیماً روی حاشیه ایمنی اثر گذاشته است. در عامل معلم، افت عملکرد محدودتر است، در حالی که عامل دانشجو در برخی مراحل حساسیت بیشتری به پیچیدگی محیط نشان می‌دهد. بنابراین، نقطه ضعف اصلی دانشجو نه ناتوانی در یادگیری مأموریت، بلکه کاهش پایداری ایمنی در شرایط دشوارتر است.

\subsection{تحلیل سپر ایمنی}
\label{subsec:shield_analysis}

شکل~\ref{fig:teacher_student_shield_dependency} رابطه میان نرخ موفقیت و نسبت مداخلات سپر را نشان می‌دهد. در عامل معلم، نرخ موفقیت پس از فاز ابتدایی در سطح بالایی باقی می‌ماند و نسبت مداخلات سپر در محدوده‌ای کنترل‌شده نوسان می‌کند. در عامل دانشجو، نرخ موفقیت همچنان قابل قبول است، اما نسبت مداخلات سپر بزرگ‌تر و نوسانی‌تر است. این تفاوت نشان می‌دهد که دانشجو برای حفظ عملکرد خود بیشتر به سپر ایمنی متکی بوده است.

\begin{figure}[htbp]
	\centering
	
	\subfloat[]{
		\begin{minipage}{0.95\textwidth}
			\centering
			
			\includegraphics[width=0.9\textwidth]
			{27_shield_dependency_decay_teacher.png}
			
			\vspace{0.35em}
			
			{\small
				\href{https://mohammadabedi1179.github.io/thesis-archive/chapters/ch6/27_shield_dependency_decay_teacher.html}
				{\textbf{مشاهده نسخه تعاملی عامل معلم با کیفیت بالاتر}}
				\hspace{0.8em}
				\raisebox{-0.4\height}{%
					\href{https://mohammadabedi1179.github.io/thesis-archive/chapters/ch6/27_shield_dependency_decay_teacher.html}{%
						\qrcode[height=1.2cm]
						{https://mohammadabedi1179.github.io/thesis-archive/chapters/ch6/27_shield_dependency_decay_teacher.html}%
					}%
				}
			}
		\end{minipage}
	}
	
	\vspace{0.7em}
	
	\subfloat[]{
		\begin{minipage}{0.95\textwidth}
			\centering
			
			\includegraphics[width=0.9\textwidth]
			{27_shield_dependency_decay_student.png}
			
			\vspace{0.35em}
			
			{\small
				\href{https://mohammadabedi1179.github.io/thesis-archive/chapters/ch6/27_shield_dependency_decay_student.html}
				{\textbf{مشاهده نسخه تعاملی عامل دانشجو با کیفیت بالاتر}}
				\hspace{0.8em}
				\raisebox{-0.4\height}{%
					\href{https://mohammadabedi1179.github.io/thesis-archive/chapters/ch6/27_shield_dependency_decay_student.html}{%
						\qrcode[height=1.2cm]
						{https://mohammadabedi1179.github.io/thesis-archive/chapters/ch6/27_shield_dependency_decay_student.html}%
					}%
				}
			}
		\end{minipage}
	}
	
	\caption{مقایسه روند نرخ موفقیت و نسبت مداخلات سپر در طول آموزش:
		الف) عامل معلم، ب) عامل دانشجو}
	\label{fig:teacher_student_shield_dependency}
\end{figure}

برای تکمیل این تحلیل، شکل~\ref{fig:teacher_student_shield_severity} شدت رخدادهای سپر را نشان می‌دهد. در عامل معلم، بخش عمده رخدادها در سطح احتیاط باقی می‌ماند و امتیاز شدت معمولاً کنترل‌شده است. در عامل دانشجو، سهم رخدادهای خطر و بحرانی در برخی بازه‌ها بیشتر است و امتیاز شدت نوسان‌های بزرگ‌تری دارد. بنابراین، وابستگی بیشتر دانشجو به سپر فقط از جنس هشدارهای سطح پایین نیست، بلکه در برخی مراحل با نزدیک‌شدن بیشتر به وضعیت‌های پرخطر همراه شده است.

\begin{figure}[htbp]
	\centering
	
	\subfloat[]{
		\begin{minipage}[t]{0.48\textwidth}
			\centering
			
			\includegraphics[width=\linewidth]
			{28_shield_severity_evolution_teacher.png}
			
			\vspace{0.35em}
			
			{\small
				\href{https://mohammadabedi1179.github.io/thesis-archive/chapters/ch6/28_shield_severity_evolution_teacher.html}
				{\textbf{مشاهده نسخه تعاملی با کیفیت بالاتر}}
				
				\par\vspace{0.3em}
				
				\href{https://mohammadabedi1179.github.io/thesis-archive/chapters/ch6/28_shield_severity_evolution_teacher.html}{%
					\qrcode[height=1.15cm]
					{https://mohammadabedi1179.github.io/thesis-archive/chapters/ch6/28_shield_severity_evolution_teacher.html}%
				}
			}
		\end{minipage}
	}
	\hfill
	\subfloat[]{
		\begin{minipage}[t]{0.48\textwidth}
			\centering
			
			\includegraphics[width=\linewidth]
			{28_shield_severity_evolution_student.png}
			
			\vspace{0.35em}
			
			{\small
				\href{https://mohammadabedi1179.github.io/thesis-archive/chapters/ch6/28_shield_severity_evolution_student.html}
				{\textbf{مشاهده نسخه تعاملی با کیفیت بالاتر}}
				
				\par\vspace{0.3em}
				
				\href{https://mohammadabedi1179.github.io/thesis-archive/chapters/ch6/28_shield_severity_evolution_student.html}{%
					\qrcode[height=1.15cm]
					{https://mohammadabedi1179.github.io/thesis-archive/chapters/ch6/28_shield_severity_evolution_student.html}%
				}
			}
		\end{minipage}
	}
	
	\caption{مقایسه ترکیب شدت رخدادهای سپر ایمنی در طول آموزش:
		الف) عامل معلم، ب) عامل دانشجو}
	\label{fig:teacher_student_shield_severity}
\end{figure}

این نتایج نشان می‌دهد که سپر ایمنی فقط یک ابزار کمکی برای جلوگیری از برخورد نیست، بلکه شاخصی مهم برای آشکارسازی کیفیت ایمنی سیاست است. عامل دانشجو از نظر موفقیت کلی به سطح قابل قبول رسیده، اما نسبت به معلم وابستگی بیشتری به سپر و رخدادهای شدیدتری در نواحی خطر دارد.

\subsection{جمع‌بندی تحلیل آموزش}
\label{subsec:training_analysis_summary}

تحلیل نتایج آموزش نشان می‌دهد که هر دو عامل از فاز ابتدایی ناپایدار عبور کرده و به رفتار هدفمند با پاداش و نرخ موفقیت مناسب رسیده‌اند. عامل معلم، به دلیل استفاده از مشاهده هندسی دقیق، در بیشتر شاخص‌ها پایدارتر است. عامل دانشجو نیز توانسته قابلیت کلی ناوبری را حفظ کند، اما در مقایسه با معلم با نوسان بیشتر، برخوردهای بیشتر، زیان منتقد بزرگ‌تر و وابستگی بالاتر به سپر ایمنی همراه است.

بنابراین، سازوکار آموزش معلم--دانشجو در انتقال سیاست از فضای هندسی به فضای ادراکی موفق بوده، اما این انتقال بدون هزینه نبوده است. هزینه اصلی در کاهش نسبی حاشیه ایمنی و افزایش حساسیت دانشجو به سناریوهای دشوار دیده می‌شود. بر همین اساس، ارزیابی نهایی سیاست دانشجو در سناریوهای مستقل از آموزش ضروری است تا مشخص شود عملکرد مشاهده‌شده در آموزش تا چه حد در شرایط آزمون نیز حفظ می‌شود.
\section{ارزیابی نهایی سیاست‌ها در سناریوی آزمون}
\label{sec:final_policy_evaluation}

پس از تحلیل روند آموزش و بررسی رفتار سیاست‌ها در طول یادگیری، در این بخش عملکرد نهایی روش پیشنهادی در سخت‌ترین سناریوی آزمون ارزیابی می‌شود. هدف این ارزیابی، صرفاً گزارش عملکرد یک سیاست منفرد نیست، بلکه مقایسه جامع میان عامل معلم، عامل دانشجو و یک خط مبنای کلاسیک مبتنی بر \lr{MPC} در شرایط مختلف مشاهده و ایمنی است. ازاین‌رو، چند حالت ارزیابی تعریف شده‌اند تا بتوان اثر نوع سیاست، منبع مشاهده و نقش سپر ایمنی را به‌صورت تفکیک‌شده بررسی کرد.

سناریوی آزمون متناظر با مرحله نهایی آموزش است. در این مرحله، محیط شامل سه مانع ایستا، چهار مانع متحرک و یک هدف متحرک است. هدف با سرعت \lr{0.5 m/s} حرکت می‌کند، محدودیت زمانی هر اپیزود برابر با \lr{100 s} در نظر گرفته شده و شعاع موفقیت نیز برابر با \lr{1.5 m} است. برای هر یک از حالت‌های ارزیابی در شبیه‌سازی، 100 اپیزود مستقل اجرا شد و شاخص‌هایی مانند پاداش اپیزود، طول اپیزود، دلیل پایان، زمان تا برخورد، حداقل فاصله از موانع، شدت و فراوانی مداخلات سپر، فرمان‌های کنترلی و شاخص‌های نرمی حرکت ثبت شدند.

\subsection{تعریف حالت‌های ارزیابی نهایی}
\label{subsec:final_eval_protocol}

برای تحلیل دقیق‌تر رفتار سامانه، شش حالت ارزیابی تعریف شد که در جدول~\ref{tab:final_eval_modes} خلاصه شده‌اند. این آزمون‌ها به‌گونه‌ای انتخاب شده‌اند که سه محور اصلی ارزیابی را پوشش دهند: نوع سیاست، کیفیت مشاهده و نقش سپر ایمنی. چهار آزمون نخست به عامل‌های یادگیرنده مربوط هستند و دو آزمون پایانی برای مقایسه با کنترل‌کننده کلاسیک \lr{MPC} در همان محیط طراحی شده‌اند.

\begin{table}[htbp]
	\centering
	\caption{حالت‌های مختلف ارزیابی نهایی در سناریوی مرحله نهم}
	\label{tab:final_eval_modes}
	\small
	\begin{tabular}{c p{3.2cm} p{3.0cm} p{3.5cm}}
		\hline
		آزمون & مدل ارزیابی‌شده & منبع مشاهده & وضعیت سپر ایمنی \\
		\hline
		آزمون اول & دانشجو & داده‌های حسگری & سپر فعال \\
		آزمون دوم & دانشجو & داده‌های حسگری & بدون مداخله فعال سپر \\
		آزمون سوم & معلم & داده‌های ایده‌آل & سپر فعال \\
		آزمون چهارم & معلم & داده‌های حسگری & سپر فعال \\
		آزمون پنجم & \lr{MPC} & داده‌های حسگری & سپر غیرفعال \\
		آزمون ششم & \lr{MPC} & داده‌های حسگری & سپر فعال \\
		\hline
	\end{tabular}
\end{table}

آزمون اول عملکرد عامل دانشجو را در حالت محافظت‌شده نشان می‌دهد؛ یعنی حالتی که تصمیم‌گیری عامل مبتنی بر داده‌های حسگری است و سپر ایمنی می‌تواند در شرایط پرخطر فرمان تولیدشده را اصلاح کند. آزمون دوم نسخه سخت‌گیرانه‌تر همین حالت است؛ در این آزمون نیز عامل دانشجو از داده‌های حسگری استفاده می‌کند، اما سپر در فرمان نهایی مداخله فعال ندارد. بنابراین، مقایسه آزمون اول و آزمون دوم نشان می‌دهد که ایمنی نهایی عامل دانشجو تا چه اندازه به لایه ایمنی وابسته است.

آزمون سوم به‌عنوان مرجع بالادستی در نظر گرفته می‌شود، زیرا در آن عامل معلم با داده‌های ایده‌آل و سپر فعال ارزیابی شده است. آزمون چهارم اثر استفاده از داده‌های حسگری را برای همین عامل، در حضور سپر فعال، نشان می‌دهد. مقایسه این دو آزمون، میزان افت عملکرد ناشی از جایگزینی مشاهده ایده‌آل با مشاهده ادراکی را مشخص می‌کند.

در نهایت، آزمون‌های پنجم و ششم برای مقایسه با یک خط مبنای غیر یادگیرنده انجام شده‌اند. در این دو آزمون، کنترل‌کننده \lr{MPC} در همان سناریوی نهایی و با داده‌های حسگری ارزیابی شده است. تفاوت این دو آزمون در وضعیت سپر ایمنی است؛ در آزمون پنجم سپر غیرفعال بوده و در آزمون ششم سپر فعال شده است. بدین ترتیب، می‌توان بررسی کرد که یک کنترل‌کننده پیش‌بین کلاسیک در همین محیط پویا و دشوار، با و بدون کمک لایه ایمنی، چه سطحی از عملکرد را ارائه می‌دهد.

\subsection{مقایسه عملکرد کلی در حالت‌های ارزیابی}
\label{subsec:final_eval_overall_comparison}

خلاصه نتایج خروجی اپیزودها در جدول~\ref{tab:final_eval_outcome_summary} ارائه شده است. این جدول نشان می‌دهد که عملکرد کلی سیاست‌ها در شرایط مختلف تفاوت معناداری دارد و حذف سپر ایمنی، کاهش کیفیت مشاهده یا استفاده از یک خط مبنای کلاسیک، مستقیماً بر نرخ موفقیت و برخورد اثر می‌گذارد.

\begin{table}[htbp]
	\centering
	\caption{خلاصه نتایج خروجی اپیزودها در شش حالت ارزیابی نهایی}
	\label{tab:final_eval_outcome_summary}
	\small
	\begin{tabular}{c p{6.8cm} c c c}
		\hline
		آزمون & حالت ارزیابی & موفقیت & اتمام زمان & برخورد \\
		\hline
		آزمون اول & دانشجو، داده حسگری، سپر فعال & 96\% & 4\% & 0\% \\
		آزمون دوم & دانشجو، داده حسگری، بدون مداخله فعال سپر & 81\% & 2\% & 17\% \\
		آزمون سوم & معلم، داده ایده‌آل، سپر فعال & 98\% & 1\% & 1\% \\
		آزمون چهارم & معلم، داده حسگری، سپر فعال & 94\% & 1\% & 5\% \\
		آزمون پنجم & \lr{MPC}، داده حسگری، سپر غیرفعال & 30\% & 0\% & 70\% \\
		آزمون ششم & \lr{MPC}، داده حسگری، سپر فعال & 50\% & 2\% & 48\% \\
		\hline
	\end{tabular}
\end{table}

در آزمون اول، عامل دانشجو با داده‌های حسگری و سپر فعال به نرخ موفقیت 96 درصد و نرخ برخورد صفر دست یافته است. این نتیجه نشان می‌دهد که ترکیب سیاست دانشجو و سپر ایمنی می‌تواند در سخت‌ترین سناریوی آزمون، عملکرد محافظت‌شده و قابل قبولی ایجاد کند. بااین‌حال، مقایسه با آزمون دوم نشان می‌دهد که این عملکرد به‌طور کامل ناشی از ایمنی ذاتی خود سیاست نیست. در آزمون دوم، با حذف مداخله فعال سپر، نرخ موفقیت به 81 درصد کاهش یافته و نرخ برخورد به 17 درصد افزایش یافته است. بنابراین، سیاست دانشجو بدون سپر فعال هنوز در شرایط دشوار، از نظر ایمنی کاملاً خودکفا نیست.

آزمون سوم بهترین عملکرد کلی را در میان همه حالات نشان می‌دهد. در این آزمون، عامل معلم با داده‌های ایده‌آل و سپر فعال به 98 درصد موفقیت و تنها 1 درصد برخورد رسیده است. این نتیجه نشان می‌دهد که در صورت دسترسی به مشاهده دقیق و کم‌خطا، سیاست معلم در سناریوی نهایی رفتار بسیار پایداری دارد. بااین‌حال، مقایسه آزمون سوم و آزمون چهارم نشان می‌دهد که جایگزینی داده‌های ایده‌آل با داده‌های حسگری موجب افت عملکرد می‌شود. در آزمون چهارم، موفقیت به 94 درصد کاهش یافته و برخورد به 5 درصد رسیده است. این اختلاف نشان می‌دهد که بخشی از محدودیت نهایی سامانه از عدم‌قطعیت مسیر ادراکی ناشی می‌شود، نه صرفاً از سیاست تصمیم‌گیری.

آزمون‌های پنجم و ششم عملکرد کنترل‌کننده \lr{MPC} را در همین سناریوی نهایی و با داده‌های حسگری نشان می‌دهند. در آزمون پنجم، یعنی \lr{MPC} بدون سپر، نرخ موفقیت تنها 30 درصد و نرخ برخورد 70 درصد است. این نتیجه نشان می‌دهد که \lr{MPC} در محیط پویا و شلوغ مرحله نهم، بدون کمک لایه ایمنی، قادر به حفظ عملکرد ایمن نیست. در آزمون ششم، با فعال‌سازی سپر ایمنی، نرخ موفقیت \lr{MPC} به 50 درصد افزایش یافته و نرخ برخورد به 48 درصد کاهش یافته است. بنابراین، سپر ایمنی توانسته عملکرد \lr{MPC} را بهبود دهد، اما حتی در حالت محافظت‌شده نیز نرخ برخورد آن به‌مراتب بیشتر از عامل‌های یادگیرنده باقی مانده است.

از دید مقایسه مستقیم، عامل دانشجو حتی در حالت بدون مداخله فعال سپر نیز عملکرد بهتری از \lr{MPC} دارد. دانشجو در آزمون دوم به 81 درصد موفقیت و 17 درصد برخورد رسیده است، در حالی که \lr{MPC} بدون سپر تنها 30 درصد موفقیت و 70 درصد برخورد نشان داده است. این اختلاف بیانگر آن است که سیاست یادگیری‌شده توانسته برای محیط پویا و دارای موانع متحرک، الگوهای تصمیم‌گیری مناسب‌تری نسبت به کنترل‌کننده پیش‌بین کلاسیک بیاموزد.

\subsection{تحلیل جامع داشبوردهای ارزیابی نهایی}
\label{subsec:final_eval_dashboard_analysis}

برای تحلیل دقیق‌تر نتایج، داشبوردهای کامل عامل‌های یادگیرنده در شکل‌های~\ref{fig:eval_dashboard_test1} تا~\ref{fig:eval_dashboard_test4} و داشبوردهای فشرده \lr{MPC} در شکل~\ref{fig:mpc_dashboards} ارائه شده‌اند. داشبوردهای عامل‌های یادگیرنده شامل پاداش اپیزودها، طول اپیزودها، ترکیب دلایل پایان، توزیع \lr{TTC}، حداقل فاصله از موانع، مداخلات سپر، فرمان‌های میانگین، شاخص تکان و نمودار همبستگی شاخص‌ها با پاداش هستند. در مقابل، داشبوردهای \lr{MPC} علاوه بر خروجی اپیزودها و هندسه ایمنی، شاخص‌های محاسباتی حل‌گر و الگوی اشباع فرمان را نیز نشان می‌دهند.

\begin{figure}[htb]
	\centering
	\includegraphics[width=0.84\textwidth]{eval_dashboard_student_active-shield.png}
	\caption{داشبورد ارزیابی نهایی برای آزمون اول: عامل دانشجو، داده‌های حسگری، سپر فعال}
	\label{fig:eval_dashboard_test1}
\end{figure}

\begin{figure}[htb]
	\centering
	\includegraphics[width=0.84\textwidth]{eval_dashboard_student_perception_passive_shield.png}
	\caption{داشبورد ارزیابی نهایی برای آزمون دوم: عامل دانشجو، داده‌های حسگری، بدون مداخله فعال سپر}
	\label{fig:eval_dashboard_test2}
\end{figure}

\begin{figure}[htb]
	\centering
	\includegraphics[width=0.84\textwidth]{eval_dashboard_teacher_use-gt_active-shield.png}
	\caption{داشبورد ارزیابی نهایی برای آزمون سوم: عامل معلم، داده‌های ایده‌آل، سپر فعال}
	\label{fig:eval_dashboard_test3}
\end{figure}

\begin{figure}[htb]
	\centering
	\includegraphics[width=0.84\textwidth]{eval_dashboard_teacher_perception_active-shield.png}
	\caption{داشبورد ارزیابی نهایی برای آزمون چهارم: عامل معلم، داده‌های حسگری، سپر فعال}
	\label{fig:eval_dashboard_test4}
\end{figure}

\begin{figure}[htbp]
	\centering
	\subfloat[]{
		\includegraphics[width=0.48\textwidth]{mpc_compact_eval_shield_off.png}
	}
	\hfill
	\subfloat[]{
		\includegraphics[width=0.48\textwidth]{mpc_compact_eval_shield_on.png}
	}
	\caption{داشبوردهای فشرده ارزیابی \lr{MPC} در سناریوی مرحله نهم: الف) آزمون پنجم، \lr{MPC} با داده‌های حسگری و سپر غیرفعال، ب) آزمون ششم، \lr{MPC} با داده‌های حسگری و سپر فعال}
	\label{fig:mpc_dashboards}
\end{figure}

در آزمون اول، بیشتر اپیزودها دارای پاداش مثبت و نسبتاً بالا هستند و تنها چند اپیزود به اتمام زمان منجر شده‌اند. پنل خروجی‌ها نشان می‌دهد که هیچ برخوردی رخ نداده است. طول اپیزودها نیز بیانگر آن است که اغلب اپیزودهای موفق پیش از رسیدن به سقف زمانی به پایان رسیده‌اند و موارد ناموفق عمدتاً از نوع نرسیدن به هدف در زمان مجاز بوده‌اند، نه برخورد. با وجود نرخ برخورد صفر، پنل‌های ایمنی نشان می‌دهند که این حالت کاملاً بدون خطر نبوده است. در چند اپیزود تعداد قابل توجهی مداخله سپر دیده می‌شود که بیانگر نقش اصلاحی واقعی سپر در پیشگیری از برخورد است. همچنین، توزیع \lr{TTC} نشان می‌دهد که اگرچه در بسیاری از اپیزودها میانگین \lr{TTC} بزرگ بوده، اما در برخی لحظات مقدار حداقل \lr{TTC} کوچک شده است؛ یعنی عامل در بعضی بازه‌ها به موقعیت‌های نزدیک به خطر وارد شده، اما سپر از تبدیل آن‌ها به شکست نهایی جلوگیری کرده است.

در آزمون دوم، تصویر متفاوتی دیده می‌شود. هرچند هنوز بخش زیادی از اپیزودها موفق هستند، اما اپیزودهای برخورد با پاداش‌های منفی به‌وضوح ظاهر شده‌اند. از آنجا که مدل ارزیابی‌شده و داده‌های ورودی همانند آزمون اول هستند و تنها تفاوت در غیرفعال بودن مداخله سپر است، افزایش نرخ برخورد در اینجا مستقیماً به نبود لایه ایمنی فعال نسبت داده می‌شود. پنل‌های حداقل فاصله از مانع و توزیع \lr{TTC} نیز نشان می‌دهند که عامل در تعداد بیشتری از اپیزودها حاشیه ایمنی خود را از دست داده است. بنابراین، مشکل آزمون دوم فقط افت بازده نیست، بلکه کاهش واقعی ایمنی ناوبری نیز مشاهده می‌شود.

آزمون سوم عملکرد مرجع عامل معلم را در بهترین شرایط مشاهده نشان می‌دهد. در این حالت، تقریباً همه اپیزودها موفق بوده‌اند و فقط تعداد بسیار محدودی به برخورد یا اتمام زمان منجر شده‌اند. این نتیجه نشان می‌دهد که وقتی عدم‌قطعیت ادراکی حذف می‌شود، سیاست معلم در سناریوی نهایی رفتار بسیار پایدار و سازگاری تولید می‌کند. بااین‌حال، حتی در این حالت نیز پنل‌های ایمنی نشان می‌دهند که سناریوی مرحله نهم ساده نیست؛ در برخی اپیزودها حداقل \lr{TTC} کاهش یافته و سپر فعال شده است. بنابراین، وجود رفتار ایمن در این آزمون صرفاً ناشی از سادگی محیط نیست، بلکه حاصل کیفیت بالای مشاهده و توانایی سیاست است.

آزمون چهارم اثر استفاده از داده‌های حسگری را برای عامل معلم نشان می‌دهد. با وجود آنکه این عامل همچنان نرخ موفقیت بالایی دارد، نسبت به آزمون سوم افت مشخصی در موفقیت و افزایش برخورد دیده می‌شود. این موضوع نشان می‌دهد که بخش مهمی از محدودیت عملکرد نهایی به کیفیت مشاهده وابسته است. پنل‌های حداقل فاصله از موانع و توزیع \lr{TTC} نیز نسبت به آزمون سوم پراکندگی بیشتری دارند و نشان می‌دهند که عامل در حضور داده‌های حسگری، بیشتر در معرض موقعیت‌های نزدیک به خطر قرار می‌گیرد. بااین‌حال، عملکرد عامل معلم در این حالت همچنان از دانشجو بدون سپر و نیز از \lr{MPC} بهتر باقی می‌ماند.

داشبوردهای \lr{MPC} در شکل~\ref{fig:mpc_dashboards} نشان می‌دهند که این کنترل‌کننده در همین سناریو و با داده‌های حسگری، با چالش جدی روبه‌رو است. در آزمون پنجم، یعنی \lr{MPC} بدون سپر، نرخ موفقیت تنها 30 درصد و نرخ برخورد 70 درصد است. نمودار هندسه ایمنی بیانگر آن است که بخش بزرگی از اپیزودهای برخورد در ناحیه فاصله کم از مانع و \lr{TTC} پایین رخ داده‌اند. همچنین، نمودار الگوی تحریک کنترل نشان می‌دهد که \lr{MPC} در بخش قابل توجهی از زمان، فرمان‌های نزدیک به اشباع تولید کرده است. این رفتار در محیطی با موانع متحرک و هدف سریع، احتمال خطا و برخورد را افزایش می‌دهد.

در آزمون ششم، با فعال‌سازی سپر ایمنی، عملکرد \lr{MPC} تا حدی بهبود یافته است. نرخ موفقیت از 30 درصد به 50 درصد افزایش یافته و نرخ برخورد از 70 درصد به 48 درصد کاهش پیدا کرده است. این نتیجه نشان می‌دهد که سپر فعال برای \lr{MPC} نیز سودمند است، اما این سودمندی برای رساندن عملکرد آن به سطح عامل‌های یادگیرنده کافی نیست. در واقع، حتی \lr{MPC} با سپر فعال نیز همچنان نرخ برخوردی به‌مراتب بالاتر از عامل دانشجو با سپر فعال و عامل معلم با داده حسگری دارد.

علاوه بر شاخص‌های خروجی، داشبوردهای \lr{MPC} اطلاعات مهمی درباره بار محاسباتی حل‌گر نیز ارائه می‌کنند. زمان حل در هر گام کنترل، در هر دو آزمون \lr{MPC} در مقایسه با گام زمانی محیط قابل توجه است و نشان می‌دهد که این رویکرد از نظر محاسباتی نیز سنگین‌تر از سیاست یادگیری‌شده است. در مقابل، عامل‌های مبتنی بر شبکه عصبی پس از آموزش، فرمان کنترل را با هزینه محاسباتی بسیار کمتری تولید می‌کنند. ازاین‌رو، برتری عامل‌های یادگیرنده در این بخش فقط به نرخ موفقیت محدود نمی‌شود، بلکه از منظر امکان استفاده عملی در سامانه بلادرنگ نیز قابل توجه است.

مقایسه کلی شش آزمون چند نتیجه مهم را روشن می‌کند. نخست آنکه بهترین عملکرد کلی به آزمون سوم، یعنی عامل معلم با داده‌های ایده‌آل و سپر فعال، تعلق دارد. دوم آنکه عامل دانشجو در آزمون اول توانسته با وجود استفاده از داده‌های حسگری، به عملکردی نزدیک به معلم برسد و در عین حال نرخ برخورد صفر را حفظ کند. سوم آنکه مقایسه آزمون اول و آزمون دوم نشان می‌دهد سپر ایمنی برای عامل دانشجو نقش حیاتی دارد، زیرا حذف آن افزایش شدید برخورد را در پی دارد. چهارم آنکه مقایسه آزمون سوم و آزمون چهارم نشان می‌دهد محدودیت ادراک یکی از گلوگاه‌های اصلی سامانه است. و در نهایت، مقایسه آزمون‌های \lr{MPC} با عامل‌های یادگیرنده نشان می‌دهد که در محیط پویا و دارای عدم‌قطعیت ادراکی، چارچوب یادگیری پیشنهادی از نظر کارایی ناوبری، ایمنی و قابلیت استفاده عملی بر خط مبنای کلاسیک \lr{MPC} برتری دارد.

در مجموع، نتایج این بخش نشان می‌دهند که چارچوب معلم--دانشجو توانسته یک سیاست ادراکی مؤثر تولید کند که در حالت محافظت‌شده، در سخت‌ترین سناریوی آزمون، عملکردی نزدیک به عامل معلم و بهتر از \lr{MPC} ارائه می‌دهد. بااین‌حال، نتایج آزمون دوم نیز به‌صراحت نشان می‌دهند که ایمنی ذاتی سیاست دانشجو هنوز کامل نیست و برای استفاده مطمئن در سناریوهای دشوار، همچنان به سپر فعال نیاز دارد.

\section{ارزیابی عملکرد اجزای سامانه عملی}
\label{sec:real_system_component_validation}

پس از ارزیابی سیاست‌ها در محیط شبیه‌سازی، پیش از بررسی اجرای کامل سیاست روی اسکوتر واقعی، لازم است عملکرد اجزای اصلی سامانه عملی به‌صورت جداگانه ارزیابی شود. دلیل این کار آن است که عملکرد نهایی سامانه تنها به کیفیت سیاست یادگیری وابسته نیست، بلکه دقت زنجیره فاصله‌سنجی، صحت تخمین وضعیت حرکتی، پایداری اندازه‌گیری حسگر اینرسی و توانایی سامانه حرکتی در اجرای فرمان‌های تولیدشده نیز بر نتیجه نهایی اثر مستقیم دارند. بنابراین، در این بخش ابتدا اجزای کلیدی سامانه عملی شامل فاصله‌سنجی استریو، حسگر اینرسی و پاسخ سامانه حرکتی بررسی می‌شوند. سپس، در بخش بعدی، اجرای سیاست آموزش‌دیده روی سامانه واقعی تحلیل خواهد شد.
\subsection{ارزیابی دقت فاصله‌سنجی استریو در حالت ایستا}
\label{subsec:stereo_depth_accuracy_test}

برای بررسی صحت زنجیره فاصله‌سنجی استریو در سامانه عملی، یک آزمون تجربی در حالت ایستا انجام شد. در این آزمون، اجسام مختلف در فاصله‌های معلوم نسبت به دوربین قرار گرفتند و فاصله واقعی آن‌ها با متر اندازه‌گیری شد. سپس در تصویر چپ تصحیح‌شده، ناحیه‌ای از جسم به‌صورت دستی انتخاب گردید و مقدار عمق از روی نقشه اختلاف مکانی همان ناحیه استخراج شد. استفاده از ناحیه به‌جای یک پیکسل منفرد باعث می‌شود اثر نویز نقطه‌ای، خطاهای موضعی نقشه اختلاف مکانی و مقادیر پرت کاهش یابد.

برای هر ناحیه انتخابی، ابتدا پیکسل‌های دارای اختلاف مکانی معتبر استخراج شدند. سپس عمق متناظر با هر پیکسل معتبر بر اساس رابطه استریو محاسبه گردید. به‌منظور کاهش اثر نقاط پرت، بخش مرکزی توزیع عمق مورد استفاده قرار گرفت و مقدار نهایی فاصله از روی مقادیر باقی‌مانده محاسبه شد. علاوه بر فاصله تخمین‌زده‌شده، شاخص‌هایی مانند خطای اندازه‌گیری، انحراف معیار عمق، درصد پیکسل‌های معتبر و کیفیت اندازه‌گیری نیز ثبت شدند.

\begin{figure}[htb]
	\centering
	\includegraphics[width=0.75\linewidth]{depth_accuracy_thesis_style.png}
	\caption{مقایسه فاصله واقعی و فاصله اندازه‌گیری‌شده در آزمون تجربی دقت فاصله‌سنجی استریو}
	\label{fig:depth_accuracy_true_measured}
\end{figure}

در مجموع، $20$ اندازه‌گیری در بازه فاصله حدود $0.48$ تا $5.20~m$ انجام شد. خلاصه نتایج این آزمون در جدول~\ref{tab:stereo_depth_accuracy_summary} ارائه شده است. میانگین خطا برابر با $0.0025~m$ به دست آمده است که نشان می‌دهد خطای تخمین عمق در این مجموعه آزمایش‌ها دارای بایاس قابل توجهی نبوده است. بااین‌حال، برای ارزیابی واقعی دقت فاصله‌سنجی، شاخص‌های قدرمطلق خطا و ریشه میانگین مربعات خطا اهمیت بیشتری دارند. در این آزمون، میانگین قدرمطلق خطا برابر با $0.0660~m$، ریشه میانگین مربعات خطا برابر با $0.0812~m$ و میانگین قدرمطلق خطای نسبی برابر با $3.50\%$ به دست آمد. این مقادیر نشان می‌دهند که زنجیره فاصله‌سنجی استریو در بازه آزمایش‌شده، برای تشکیل مشاهده عامل و تخمین فاصله نسبی هدف و موانع قابل استفاده است.

\begin{table}[htbp]
	\centering
	\caption{خلاصه نتایج آزمون تجربی دقت فاصله‌سنجی استریو}
	\label{tab:stereo_depth_accuracy_summary}
	\renewcommand{\arraystretch}{1.35}
	\begin{tabular}{lc}
		\hline
		\textbf{شاخص ارزیابی} & \textbf{مقدار} \\
		\hline
		تعداد اندازه‌گیری‌ها & $20$ \\
		بازه فاصله واقعی & $0.48$ تا $5.20~m$ \\
		میانگین خطا & $0.0025~m$ \\
		میانگین قدرمطلق خطا & $0.0660~m$ \\
		ریشه میانگین مربعات خطا & $0.0812~m$ \\
		میانگین قدرمطلق خطای نسبی & $3.50\%$ \\
		میانگین پوشش پیکسل‌های معتبر & $95.95\%$ \\
		\hline
	\end{tabular}
\end{table}

شکل~\ref{fig:depth_accuracy_sample} نمونه‌ای از روند آزمون را نشان می‌دهد. در این شکل، تصویر چپ تصحیح‌شده، ناحیه انتخابی روی جسم و نقشه اختلاف مکانی متناظر نمایش داده شده است. مقایسه این خروجی‌ها نشان می‌دهد که نواحی دارای بافت مناسب و مرزهای قابل تفکیک، خروجی عمق پایدارتر و قابل اعتمادتر تولید می‌کنند. در مقابل، نواحی کم‌بافت، سطوح بازتابنده یا مرزهای ضعیف می‌توانند باعث کاهش کیفیت تخمین اختلاف مکانی شوند.

\begin{figure}[htbp]
	\centering
	\includegraphics[width=0.95\linewidth]{depth_accuracy_sample.png}
	\caption{نمونه‌ای از آزمون تجربی فاصله‌سنجی استریو با انتخاب ناحیه روی جسم و استخراج عمق از نقشه اختلاف مکانی}
	\label{fig:depth_accuracy_sample}
\end{figure}

بر اساس این آزمون، در زنجیره ادراک سامانه عملی از تخمین عمق مبتنی بر ناحیه استفاده شد. به این ترتیب، فاصله هدف و موانع از ناحیه معتبر داخل کادر تشخیص یا نشانگر هدف استخراج می‌شود و از اتکا به مقدار عمق یک پیکسل منفرد جلوگیری می‌گردد. با وجود دقت مناسب در حالت ایستا، باید توجه داشت که این آزمون اثراتی مانند حرکت ربات، لرزش دوربین، تغییرات سریع نور، تأخیر پردازش و خطاهای ناشی از تشخیص هدف یا مانع را به‌طور کامل پوشش نمی‌دهد. بنابراین، نتایج این بخش نشان‌دهنده دقت پایه زنجیره فاصله‌سنجی استریو هستند و در ادامه، عملکرد کل سامانه در اجرای واقعی نیز به‌صورت جداگانه بررسی می‌شود.
\subsection{ارزیابی زنجیره تخمین فاصله نسبی در اجرای واقعی اسکوتر}
\label{subsec:real_scooter_relative_depth_accuracy}

پس از ارزیابی دقت پایه فاصله‌سنجی استریو در حالت ایستا، عملکرد زنجیره تخمین فاصله نسبی در شرایط واقعی‌تر و هنگام حرکت اسکوتر نیز بررسی شد. در این آزمون، هدف و موانع در محیط ثابت نگه داشته شدند، اما اسکوتر در طول مسیر حرکت کرد. موقعیت اولیه اسکوتر و موقعیت هدف و موانع به برنامه داده شد و فاصله واقعی نسبی در هر ثانیه، بر اساس موقعیت تخمینی اسکوتر محاسبه گردید. موقعیت اسکوتر نیز از ترکیب اطلاعات حسگر اینرسی و انکودر چرخ‌ها به دست آمد. بنابراین، مقدار مرجع مورد استفاده در این آزمون یک اندازه‌گیری مستقیم مستقل در هر لحظه نیست، بلکه برآوردی مبتنی بر تخمین حرکت اسکوتر است.

به همین دلیل، خطای گزارش‌شده در این آزمون نباید به‌عنوان خطای خالص دوربین استریو تفسیر شود. این خطا حاصل ترکیب چند منبع خطا است؛ از جمله خطای نقشه اختلاف مکانی، خطای تشخیص ناحیه هدف یا مانع، لرزش و حرکت دوربین، خطای تخمین موقعیت اسکوتر از حسگر اینرسی و انکودرها، خطای هم‌زمانی داده‌ها و خطای تبدیل مختصات بین قاب دوربین، بدنه اسکوتر و محیط. بااین‌حال، این آزمون از نظر کاربردی اهمیت زیادی دارد، زیرا نشان می‌دهد زنجیره ادراک و تخمین فاصله در شرایطی که برای اجرای سیاست واقعی استفاده می‌شود، چه سطحی از خطا را تولید می‌کند.

\begin{figure}[htbp]
	\centering
	\includegraphics[width=\textwidth]{00_combined_first_four_dashboard.png}
	\caption{تحلیل تخمین فاصله نسبی هدف و مانع در اجرای واقعی اسکوتر؛ فاصله مرجع در هر ثانیه از موقعیت تخمینی اسکوتر بر اساس ترکیب حسگر اینرسی و انکودرها محاسبه شده است}
	\label{fig:real_scooter_depth_dashboard}
\end{figure}

مطابق شکل~\ref{fig:real_scooter_depth_dashboard}، خطای تخمین فاصله در اجرای واقعی اسکوتر به‌طور محسوسی بیشتر از آزمون ایستای فاصله‌سنجی استریو است. برای نشانگر هدف، مقدار \lr{RMSE} داده خام برابر با حدود $0.55~m$ و برای موانع برابر با حدود $0.65~m$ است. این افزایش خطا قابل انتظار است، زیرا در این آزمون علاوه بر خطای خود استریو، خطاهای ناشی از حرکت اسکوتر، لرزش دوربین، تغییر زاویه دید، تأخیر پردازش و خطای تخمین موقعیت اسکوتر نیز وارد محاسبه می‌شوند.

نمودارهای مقایسه فاصله تخمینی و فاصله مرجع نشان می‌دهند که با وجود افزایش پراکندگی نسبت به حالت ایستا، روند کلی تغییر فاصله در بسیاری از نمونه‌ها حفظ شده است. بااین‌حال، فاصله‌های تخمینی در برخی نواحی دارای بیش‌برآورد یا کم‌برآورد قابل توجه هستند. این رفتار نشان می‌دهد که خطاها صرفاً نویز تصادفی نیستند، بلکه می‌توانند از خطاهای سیستماتیک در تخمین موقعیت، کالیبراسیون، تبدیل مختصات یا کیفیت نقشه اختلاف مکانی نیز تأثیر بگیرند.

نمودارهای خطا بر حسب سرعت خطی و سرعت زاویه‌ای نشان می‌دهند که پراکندگی خطا در شرایط حرکتی افزایش می‌یابد. به‌ویژه در سرعت‌های زاویه‌ای بالاتر، احتمال افزایش خطا بیشتر است، زیرا چرخش اسکوتر باعث تغییر سریع میدان دید، ناپایداری ناحیه تشخیص، تاری حرکتی و افزایش حساسیت نسبت به خطای هم‌زمانی میان داده‌های ادراکی و داده‌های حرکتی می‌شود. بنابراین، بخشی از افت دقت در اجرای واقعی به ماهیت دینامیکی سامانه و نه فقط به کیفیت الگوریتم استریو مربوط است.

همچنین، خطای مربوط به موانع نسبت به نشانگر هدف بیشتر است. این تفاوت قابل انتظار است، زیرا نشانگر هدف دارای ظاهر مشخص‌تر و ناحیه تشخیص پایدارتر است، در حالی که موانع ممکن است بافت کمتر، مرز ضعیف‌تر، زاویه دید نامناسب‌تر یا تشخیص ناپایدارتر داشته باشند. در نتیجه، استخراج عمق برای مانع در شرایط واقعی دشوارتر از نشانگر هدف است.

نتایج این آزمون نشان می‌دهند که اگرچه دقت تخمین فاصله در اجرای واقعی نسبت به آزمون ایستا کاهش می‌یابد، خروجی زنجیره ادراک همچنان روند کلی فاصله نسبی را حفظ می‌کند و برای تشکیل مشاهده عامل قابل استفاده است. بااین‌حال، خطاهای مشاهده‌شده ضرورت استفاده از فیلترهای پایدارساز، سپر ایمنی و سیاست مقاوم در برابر عدم‌قطعیت ادراکی را توجیه می‌کنند. بنابراین، این آزمون نقش واسطی میان ارزیابی آزمایشگاهی فاصله‌سنجی استریو و اجرای کامل سیاست روی اسکوتر واقعی دارد.

\subsection{ارزیابی خروجی جهت‌گیری حسگر اینرسی}
\label{subsec:imu_yaw_evaluation}

برای بررسی رفتار خروجی‌های جهت‌گیری حسگر اینرسی و انتخاب سیگنال مناسب برای تخمین زاویه سر اسکوتر، یک آزمون چرخش مرجع انجام شد. در این آزمون، حسگر در زاویه‌های مرجع مختلف شامل $0^\circ$، $90^\circ$، $180^\circ$، $270^\circ$ و $360^\circ$ قرار گرفت و خروجی‌های مختلف جهت‌گیری ثبت شدند. هدف این آزمون، مقایسه رفتار نسبی خروجی‌ها در برابر چرخش‌های معلوم و انتخاب خروجی پایدارتر برای استفاده در اودومتری سامانه بود.

همان‌طور که در شکل~\ref{fig:imu_yaw_heading_test} نشان داده شده است، خروجی \lr{Euler yaw} روند تغییر زاویه را در بازه‌های مرجع به‌صورت پیوسته‌تر و قابل تفسیرتر دنبال می‌کند. در مقابل، برخی خروجی‌های دیگر ممکن است به دلیل نحوه نمایش زاویه، تغییر علامت، ناپیوستگی در اطراف مرزهای زاویه‌ای یا نیاز به تبدیل از کواترنیون، رفتار مستقیم و ساده‌ای برای استفاده در اودومتری صفحه‌ای نداشته باشند. از آنجا که حرکت اسکوتر در این پژوهش عمدتاً روی صفحه انجام می‌شود، زاویه \lr{yaw} مهم‌ترین مؤلفه جهت‌گیری برای تخمین وضعیت حرکتی است.

بر همین اساس، در سامانه عملی از خروجی \lr{Euler yaw} پس از صفرسازی اولیه و بازگشایی زاویه استفاده شد. صفرسازی اولیه باعث می‌شود جهت اولیه اسکوتر به‌عنوان مرجع محلی در نظر گرفته شود و بازگشایی زاویه نیز از ایجاد ناپیوستگی‌های مصنوعی در عبور از مرزهای زاویه‌ای جلوگیری می‌کند. بنابراین، خروجی انتخاب‌شده برای ترکیب با اطلاعات انکودر چرخ‌ها و تخمین جهت حرکت اسکوتر مورد استفاده قرار گرفت.

\begin{figure}[htb]
	\centering
	\includegraphics[width=0.95\linewidth]{desk_quat_test_yaw_comparison.png}
	\caption{مقایسه خروجی‌های جهت‌گیری حسگر \lr{IMU} در آزمون چرخش‌های مرجع و انتخاب خروجی \lr{Euler yaw} برای تخمین جهت‌گیری اسکوتر}
	\label{fig:imu_yaw_heading_test}
\end{figure}

\subsection{ارزیابی پاسخ حلقه‌بسته موتورهای محرک}
\label{subsec:motor_rpm_response_test}

برای بررسی توانایی سامانه حرکتی در اجرای فرمان‌های سرعت، پاسخ چهار موتور محرک در برابر فرمان‌های پله‌ای سرعت دورانی ارزیابی شد. در این آزمون، برای هر موتور به‌صورت جداگانه مجموعه‌ای از مقادیر مرجع شامل $5$، $10$، $15$، $20$ و $30~RPM$ اعمال شد و سرعت اندازه‌گیری‌شده موتور و خروجی \lr{PWM} متناظر ثبت گردید. در هر اجرا، تنها یک موتور فعال بود و خروجی سایر موتورها صفر نگه داشته شد تا پاسخ مستقل هر موتور بررسی شود.

در سامانه عملی، سرعت هر موتور بر اساس پالس‌های حسگر هال اندازه‌گیری می‌شود و کنترل‌کننده سرعت از ترکیب مؤلفه پیش‌خور و کنترل‌کننده \lr{PI} استفاده می‌کند. علاوه بر این، برای کاهش اثر ناحیه مرده، اصطکاک سکون و تغییرات ناگهانی فرمان، جبران‌سازی ناحیه مرده، محدودیت نرخ تغییر فرمان و سازوکار ضد اشباع نیز در نظر گرفته شده است. ازاین‌رو، نتایج این بخش نباید به‌عنوان پاسخ خام موتور تفسیر شوند، بلکه پاسخ حلقه‌بسته مجموعه موتور، درایور، حسگر و کنترل‌کننده را نشان می‌دهند.

\begin{figure}[p]
	\centering
	\includegraphics[width=0.98\textwidth]{motor_sweep_overlay_all.png}
	\caption{مقایسه پاسخ حلقه‌بسته چهار موتور محرک در آزمون جاروب سرعت. در این شکل، پاسخ سرعت اندازه‌گیری‌شده، خروجی \lr{PWM}، پاسخ هر موتور در سطوح مختلف مرجع، و توزیع مقادیر \lr{PWM} برای چهار موتور به‌صورت تجمیعی نمایش داده شده است.}
	\label{fig:motor_sweep_overlay}
\end{figure}

شکل~\ref{fig:motor_sweep_overlay} مقایسه تجمیعی پاسخ چهار موتور را نشان می‌دهد. در ردیف‌های بالایی، روند کلی سرعت اندازه‌گیری‌شده و خروجی \lr{PWM} در طول کل آزمون برای چهار موتور به‌صورت هم‌پوشان نمایش داده شده است. در ردیف‌های بعدی نیز پاسخ سرعت، پاسخ \lr{PWM} و توزیع مقادیر \lr{PWM} برای هر سطح مرجع به‌صورت مجزا مقایسه شده‌اند. این نمایش تجمیعی امکان مقایسه مستقیم عملکرد موتورها را فراهم می‌کند.

نتایج نشان می‌دهند که در فرمان‌های پایین، به‌ویژه در $5~RPM$، پاسخ موتورها ناپایدارتر و کمتر خطی است. این رفتار قابل انتظار است، زیرا در سرعت‌های پایین اثر اصطکاک سکون، ناحیه مرده درایور، لقی مکانیکی و کوانتیزه‌شدن اندازه‌گیری پالس‌ها بیشتر آشکار می‌شود. در این ناحیه، بخشی از فرمان \lr{PWM} صرف غلبه بر اصطکاک اولیه می‌شود و تا پیش از عبور از یک آستانه مشخص، حرکت پایدار ایجاد نمی‌گردد.

در $10~RPM$، دنبال‌کردن مرجع بهتر می‌شود، اما هنوز اختلاف محسوسی میان موتورها وجود دارد. در سرعت‌های میانی، یعنی $15$ و $20~RPM$، تفاوت در رفتار چهار موتور واضح‌تر می‌شود. برخی موتورها با مقدار \lr{PWM} کمتر به سرعت مرجع نزدیک می‌شوند، در حالی که برخی دیگر برای رسیدن به همان سرعت به \lr{PWM} بیشتری نیاز دارند یا نوسان و بیش‌جهش بیشتری نشان می‌دهند. در این میان، یکی از موتورها نسبت به سایرین پاسخ کندتر و نیاز بیشتر به \lr{PWM} را نشان می‌دهد که بیانگر ناهمسانی عملگرها در سامانه واقعی است.

در $30~RPM$ هر چهار موتور به‌طور کلی وارد ناحیه عملکرد باثبات‌تری می‌شوند و سرعت مرجع را بهتر دنبال می‌کنند. بااین‌حال، حتی در این سطح نیز اختلاف میان مقدار \lr{PWM} لازم برای حفظ سرعت و میزان نوسان خروجی میان موتورها باقی می‌ماند. این موضوع نشان می‌دهد که تقارن کامل میان چهار موتور وجود ندارد و پاسخ نهایی سامانه حرکتی از تفاوت‌های ذاتی عملگرها تأثیر می‌پذیرد.

نمودارهای \lr{PWM} و هیستوگرام‌های متناظر نیز وجود ناحیه مرده و تفاوت میان راه‌اندازی و حرکت پایدار را به‌خوبی نشان می‌دهند. در بسیاری از موارد، برای شروع حرکت مقدار \lr{PWM} بیشتری لازم است، اما پس از راه‌افتادن موتور، مقدار \lr{PWM} موردنیاز برای حفظ سرعت کاهش می‌یابد. این رفتار با تفاوت میان اصطکاک سکون و اصطکاک حرکتی سازگار است. همچنین، در برخی سطوح مرجع، پس از رسیدن سرعت به مقدار مطلوب، کنترل‌کننده با کاهش تدریجی \lr{PWM} تلاش می‌کند اثر بیش‌جهش را محدود کند.

در مجموع، نتایج این آزمون نشان می‌دهند که سامانه محرک برای اجرای فرمان‌های سرعت در سامانه عملی قابل استفاده است، اما پاسخ آن ایده‌آل نیست و با ناحیه مرده، تأخیر، بیش‌جهش و ناهمسانی میان موتورها همراه است. این نتیجه در تفسیر آزمایش‌های عملی ناوبری اهمیت زیادی دارد، زیرا بخشی از اختلاف میان رفتار واقعی اسکوتر و رفتار مورد انتظار در شبیه‌سازی می‌تواند ناشی از همین محدودیت‌های عملگرها و کنترل‌کننده سرعت باشد.
\subsection{ارزیابی اجرای سیاست روی سامانه واقعی}
\label{subsec:real_world_policy_validation}

پس از ارزیابی اجزای اصلی سامانه عملی، در این بخش عملکرد سیاست نهایی در اجرای واقعی روی اسکوتر بررسی می‌شود. هدف این آزمون، ارزیابی اولیه قابلیت انتقال سیاست از محیط شبیه‌سازی به سامانه فیزیکی و سنجش رفتار زنجیره کامل ادراک، تصمیم‌گیری و فرمان‌دهی در شرایط واقعی است. به‌دلیل محدودیت‌های زمانی، ایمنی و اجرایی، این ارزیابی در قالب 20 اپیزود مستقل انجام شد. بنابراین، نتایج این بخش باید به‌عنوان یک اعتبارسنجی عملی اولیه تفسیر شوند، نه یک ارزیابی آماری گسترده در مقیاس آزمون‌های شبیه‌سازی.

در هر اپیزود، سیاست روی اسکوتر واقعی اجرا شد و ورودی‌های عامل از زنجیره حسگری سامانه شامل دوربین استریو، حسگر اینرسی و انکودر چرخ‌ها تأمین گردید. خروجی هر اپیزود از نظر پاداش نهایی، طول اپیزود، دلیل پایان و شاخص زمان تا برخورد ثبت شد. شکل~\ref{fig:real_world_validation_results} خلاصه نتایج این آزمون‌ها را نشان می‌دهد.

\begin{figure}[htbp]
	\centering
	\includegraphics[width=0.92\linewidth]{real_world_validation_dashboard.png}
	\caption{خلاصه نتایج اعتبارسنجی عملی سیاست روی اسکوتر واقعی شامل پاداش اپیزودها، طول اپیزودها، ترکیب دلایل پایان و توزیع زمان تا برخورد}
	\label{fig:real_world_validation_results}
\end{figure}

خلاصه نتایج خروجی اپیزودها در جدول~\ref{tab:real_world_validation_summary} ارائه شده است. از مجموع 20 اپیزود، 16 اپیزود با موفقیت به هدف رسیده‌اند، 3 اپیزود به اتمام زمان منتهی شده‌اند و تنها 1 اپیزود با برخورد خاتمه یافته است. به این ترتیب، نرخ موفقیت برابر با 80\%، نرخ اتمام زمان برابر با 15\% و نرخ برخورد برابر با 5\% به دست آمده است.

\begin{table}[htbp]
	\centering
	\caption{خلاصه نتایج اجرای سیاست روی اسکوتر واقعی}
	\label{tab:real_world_validation_summary}
	\renewcommand{\arraystretch}{1.3}
	\begin{tabular}{|c|c|c|}
		\hline
		\textbf{نوع پایان اپیزود} & \textbf{تعداد} & \textbf{درصد} \\
		\hline
		رسیدن به هدف & 16 & 80\% \\
		\hline
		اتمام زمان & 3 & 15\% \\
		\hline
		برخورد & 1 & 5\% \\
		\hline
		مجموع & 20 & 100\% \\
		\hline
	\end{tabular}
\end{table}

مطابق شکل~\ref{fig:real_world_validation_results}، بیشتر اپیزودهای موفق دارای پاداش مثبت و نسبتاً بالا هستند و مقدار پاداش آن‌ها عمدتاً در بازه حدود 70 تا 115 قرار گرفته است. این موضوع نشان می‌دهد که سیاست در بخش عمده‌ای از اجراهای واقعی توانسته است مأموریت ناوبری را با عملکرد قابل قبول به پایان برساند. در مقابل، اپیزودهای ناموفق با پاداش‌های منفی مشخص شده‌اند. سه اپیزود منتهی به اتمام زمان دارای پاداش‌های منفی متوسط هستند، در حالی که تنها اپیزود برخورد، کمترین پاداش را تولید کرده است. میانگین کل پاداش نیز برابر با حدود 52.6 به دست آمده که با توجه به وجود چند اپیزود ناموفق، مقدار قابل قبولی برای اجرای واقعی محسوب می‌شود.

نمودار طول اپیزودها نشان می‌دهد که سه اپیزود منتهی به اتمام زمان دقیقاً در سقف طول مجاز اپیزود قرار گرفته‌اند که نشان‌دهنده نرسیدن عامل به هدف در بازه زمانی تعریف‌شده است. در مقابل، بیشتر اپیزودهای موفق در طول‌هایی به‌مراتب کمتر از حد نهایی پایان یافته‌اند. طول اپیزودهای موفق تقریباً در بازه 50 تا 290 گام قرار دارد. این الگو نشان می‌دهد که شکست‌های مشاهده‌شده در آزمون عملی بیشتر از جنس ناتوانی در تکمیل مأموریت در زمان محدود بوده‌اند و نه بروز مکرر برخورد.

نمودار دایره‌ای نتایج نیز همین جمع‌بندی را تأیید می‌کند. سهم غالب خروجی‌ها مربوط به اپیزودهای موفق است و نرخ برخورد بسیار پایین باقی مانده است. از این نظر، سامانه عملی توانسته است در اغلب اجراها رفتار قابل قبولی از خود نشان دهد، هرچند هنوز پایداری آن به سطحی نرسیده است که همه اپیزودها را با موفقیت به پایان برساند.

توزیع \lr{TTC} نیز نکات مهمی را در مورد رفتار ایمنی سامانه نشان می‌دهد. همان‌طور که در شکل~\ref{fig:real_world_validation_results} مشاهده می‌شود، در بسیاری از اپیزودها حداقل \lr{TTC} در مقادیر کوچک متمرکز است، در حالی که میانگین \lr{TTC} در تعداد زیادی از اپیزودها در بازه‌های بزرگ‌تر، حدود 35 تا 50 ثانیه، قرار دارد. این الگو نشان می‌دهد که اگرچه سامانه در طول بخش بزرگی از هر اپیزود فاصله زمانی مناسبی تا برخورد حفظ کرده است، اما در برخی لحظات گذرا به وضعیت‌های نزدیک به خطر نیز وارد شده است. بنابراین، رسیدن به هدف در اجرای واقعی الزاماً به معنای حرکت با حاشیه ایمنی یکنواخت و بزرگ در تمام طول مسیر نیست.

در مجموع، نتایج این آزمون نشان می‌دهد که سیاست آموزش‌دیده توانسته است تا حد قابل قبولی از شبیه‌سازی به سامانه واقعی منتقل شود. دستیابی به نرخ موفقیت 80\% در کنار نرخ برخورد 5\% بیانگر آن است که زنجیره کامل سامانه، شامل ادراک، تخمین وضعیت، تصمیم‌گیری و اعمال فرمان، در محیط واقعی کارکرد عملی داشته است. بااین‌حال، وجود سه اپیزود منتهی به اتمام زمان و یک اپیزود برخورد نشان می‌دهد که شکاف میان شبیه‌سازی و واقعیت همچنان به‌طور کامل از بین نرفته است و عملکرد سامانه در شرایط واقعی همچنان تحت تأثیر خطاهای ادراکی، عدم‌قطعیت تخمین وضعیت، تغییرات محیطی و محدودیت‌های دینامیکی پلتفرم قرار دارد.

بنابراین، این نتایج را می‌توان به‌عنوان یک اعتبارسنجی عملی موفق اما اولیه در نظر گرفت. سامانه توانسته است در اکثر اجراها رفتار هدفمند و ایمن از خود نشان دهد، اما برای دستیابی به پایداری بالاتر در محیط واقعی، هنوز به بهبود بیشتر در بخش ادراک، تخمین حالت و تنظیم سیاست کنترلی نیاز است.

\section{بحث کلی نتایج}
\label{sec:chapter6_general_discussion}

نتایج ارائه‌شده در این فصل نشان می‌دهند که ارزیابی سامانه پیشنهادی باید در سه سطح تفسیر شود: نخست، عملکرد سیاست‌ها در محیط شبیه‌سازی؛ دوم، کیفیت اجزای ادراکی و حرکتی سامانه عملی؛ و سوم، رفتار کل سامانه در اجرای واقعی روی اسکوتر. این تفکیک اهمیت دارد، زیرا عملکرد نهایی سامانه تنها به کیفیت سیاست یادگیری وابسته نیست و خطاهای ادراک، تخمین وضعیت، پاسخ غیرایده‌آل موتورها و محدودیت‌های فیزیکی پلتفرم نیز بر نتیجه نهایی اثر مستقیم دارند.

در سطح شبیه‌سازی، نتایج آموزش نشان دادند که عامل معلم و عامل دانشجو هر دو توانسته‌اند الگوی کلی ناوبری به سمت هدف و اجتناب از موانع را بیاموزند. عامل معلم به‌دلیل استفاده از مشاهده ایده‌آل، عملکرد پایدارتر و کم‌نوسان‌تری دارد، در حالی که عامل دانشجو با داده‌های حسگری با عدم‌قطعیت بیشتری مواجه است. بااین‌حال، نتایج آزمون نهایی نشان دادند که عامل دانشجو در حالت محافظت‌شده با سپر فعال توانسته است به عملکردی نزدیک به عامل معلم برسد. دستیابی به نرخ موفقیت بالا و برخورد صفر در آزمون دانشجو با سپر فعال نشان می‌دهد که چارچوب معلم--دانشجو توانسته یک سیاست ادراکی قابل استفاده برای ناوبری تولید کند.

با وجود این، مقایسه آزمون دانشجو با سپر فعال و بدون مداخله فعال سپر نشان داد که ایمنی ذاتی سیاست دانشجو هنوز کامل نیست. حذف مداخله فعال سپر باعث کاهش نرخ موفقیت و افزایش برخورد شد. بنابراین، نتیجه اصلی این نیست که سیاست دانشجو به‌تنهایی کاملاً ایمن است، بلکه نتیجه دقیق‌تر آن است که ترکیب سیاست دانشجو و سپر ایمنی می‌تواند در سناریوی دشوار مرحله نهم عملکرد محافظت‌شده قابل قبولی ایجاد کند. این موضوع نشان می‌دهد که سپر ایمنی در سامانه پیشنهادی نقش تزئینی ندارد، بلکه یکی از اجزای اصلی برای تبدیل سیاست یادگیری‌شده به یک سامانه قابل استفاده در شرایط دشوار است.

مقایسه با کنترل‌کننده \lr{MPC} نیز اهمیت روش پیشنهادی را روشن‌تر می‌کند. نتایج نشان دادند که \lr{MPC} در همان سناریوی نهایی و با داده‌های حسگری، چه با سپر غیرفعال و چه با سپر فعال، عملکرد ضعیف‌تری نسبت به عامل‌های یادگیرنده دارد. نرخ برخورد بالای \lr{MPC} نشان می‌دهد که در محیطی شامل موانع متحرک، هدف متحرک و مشاهده حسگری نامطمئن، کنترل‌کننده پیش‌بین با مدل ساده‌شده و افق محدود نمی‌تواند به‌تنهایی رفتار ایمن و پایدار تولید کند. همچنین، هزینه محاسباتی بالاتر \lr{MPC} نسبت به سیاست عصبی، استفاده بلادرنگ از آن را دشوارتر می‌سازد. از این نظر، سیاست یادگیری‌شده پس از آموزش، از نظر سرعت تولید فرمان و سازگاری با شرایط پیچیده محیطی مزیت عملی دارد.

در سطح سامانه عملی، نتایج ارزیابی اجزا نشان دادند که زنجیره فاصله‌سنجی استریو در حالت ایستا دقت مناسبی دارد، اما در اجرای واقعی و هنگام حرکت اسکوتر، خطای تخمین فاصله افزایش می‌یابد. این افزایش خطا فقط به خود دوربین استریو مربوط نیست، بلکه از ترکیب خطای ادراک، لرزش دوربین، حرکت پلتفرم، تأخیر پردازش، خطای تخمین موقعیت از حسگر اینرسی و انکودرها، و خطای تبدیل مختصات ناشی می‌شود. بنابراین، شکاف میان عملکرد شبیه‌سازی و اجرای واقعی تا حد زیادی قابل انتظار است و باید در تحلیل نتایج نهایی لحاظ شود.

ارزیابی حسگر اینرسی و پاسخ حلقه‌بسته موتورها نیز نشان داد که سامانه عملی با یک پلتفرم ایده‌آل فاصله دارد. خروجی \lr{Euler yaw} پس از صفرسازی و بازگشایی زاویه برای تخمین جهت‌گیری مناسب‌تر تشخیص داده شد، اما این تخمین همچنان در زنجیره اودومتری با خطای انکودرها و حرکت واقعی ترکیب می‌شود. از سوی دیگر، پاسخ موتورها در سرعت‌های پایین تحت تأثیر ناحیه مرده، اصطکاک سکون و ناهمسانی میان موتورها قرار دارد. بنابراین، بخشی از خطاهای مشاهده‌شده در اجرای واقعی سیاست می‌تواند ناشی از تفاوت پاسخ عملگرهای واقعی با فرضیات ساده‌شده محیط شبیه‌سازی باشد.

در نهایت، آزمون‌های واقعی روی اسکوتر نشان دادند که سامانه پیشنهادی توانسته است در بیشتر اپیزودها مأموریت را با موفقیت انجام دهد. دستیابی به 16 موفقیت از 20 اپیزود، همراه با 3 اتمام زمان و تنها 1 برخورد، نشان می‌دهد که زنجیره کامل ادراک، تخمین وضعیت، تصمیم‌گیری و اعمال فرمان روی سامانه واقعی قابل اجرا بوده است. بااین‌حال، این نتایج باید به‌عنوان اعتبارسنجی عملی اولیه تفسیر شوند، نه ارزیابی آماری گسترده. تعداد محدود اپیزودهای واقعی، محدودیت‌های ایمنی آزمایش، و وابستگی نتیجه به شرایط محیطی باعث می‌شود که تعمیم این نتایج به همه شرایط عملی با احتیاط انجام شود.

به‌طور کلی، نتایج فصل نشان می‌دهند که روش پیشنهادی از نظر یادگیری سیاست ناوبری، انتقال به مشاهده حسگری و اجرای اولیه روی سامانه واقعی موفق بوده است. در عین حال، نتایج به‌روشنی نشان می‌دهند که برای استفاده پایدارتر و ایمن‌تر در محیط واقعی، بهبود دقت ادراک، کاهش تأخیر پردازش، تنظیم دقیق‌تر سامانه حرکتی، افزایش مقاومت سیاست نسبت به خطاهای حسگری و حفظ لایه ایمنی فعال همچنان ضروری است.

\section{جمع‌بندی فصل}
\label{sec:chapter6_summary}

در این فصل، عملکرد سامانه پیشنهادی از مرحله آموزش تا اجرای عملی مورد ارزیابی قرار گرفت. ابتدا روند آموزش عامل معلم و عامل دانشجو در محیط شبیه‌سازی تحلیل شد. نتایج نشان دادند که هر دو عامل توانسته‌اند رفتار ناوبری به سمت هدف و اجتناب از موانع را بیاموزند، اما عامل دانشجو به‌دلیل استفاده از داده‌های حسگری با نوسان و عدم‌قطعیت بیشتری نسبت به عامل معلم مواجه است.

در ادامه، سیاست‌های نهایی در سخت‌ترین سناریوی آزمون، شامل سه مانع ایستا، چهار مانع متحرک و هدف متحرک، ارزیابی شدند. نتایج نشان دادند که عامل دانشجو با داده‌های حسگری و سپر فعال به نرخ موفقیت بالا و برخورد صفر دست یافته است. بااین‌حال، حذف مداخله فعال سپر باعث افزایش برخورد شد و نشان داد که ایمنی ذاتی سیاست دانشجو هنوز کامل نیست. همچنین، مقایسه با \lr{MPC} نشان داد که روش یادگیری پیشنهادی در این محیط پویا و دارای عدم‌قطعیت ادراکی، عملکرد بهتری نسبت به خط مبنای کلاسیک دارد.

سپس اجزای اصلی سامانه عملی شامل زنجیره فاصله‌سنجی استریو، تخمین فاصله نسبی در حرکت واقعی، خروجی جهت‌گیری حسگر اینرسی و پاسخ حلقه‌بسته موتورهای محرک بررسی شدند. نتایج این بخش نشان دادند که هرچند اجزای سامانه برای اجرای عملی قابل استفاده‌اند، اما خطاهای ادراکی، ناحیه مرده موتورها، ناهمسانی عملگرها و خطای تخمین وضعیت از عوامل مهم اختلاف میان شبیه‌سازی و اجرای واقعی هستند.

در پایان، سیاست نهایی روی اسکوتر واقعی در قالب 20 اپیزود مستقل اجرا شد. نتایج این آزمون شامل 16 موفقیت، 3 اتمام زمان و 1 برخورد بود. این نتیجه نشان داد که زنجیره کامل سامانه در محیط واقعی قابل اجراست و سیاست آموزش‌دیده توانسته است در اغلب اجراها رفتار هدفمند و قابل قبول تولید کند. بااین‌حال، وجود اپیزودهای ناموفق نشان می‌دهد که انتقال از شبیه‌سازی به واقعیت هنوز با چالش‌هایی همراه است.

در مجموع، این فصل نشان داد که چارچوب پیشنهادی از نظر آموزش، ارزیابی شبیه‌سازی و اجرای اولیه واقعی عملکرد قابل قبولی دارد. نتیجه اصلی آن است که ترکیب سیاست دانشجو، داده‌های حسگری و سپر ایمنی می‌تواند یک راهکار عملی برای ناوبری اسکوتر خودران در محیط‌های پویا فراهم کند، اما استفاده مطمئن از آن در شرایط واقعی همچنان نیازمند بهبود ادراک، کنترل حرکتی و سازوکارهای ایمنی است.